% !TeX program = xelatex
% ECONOMIA: LEGGERE E CAPIRE - versione di lavoro, settembre 2026
% FILE AUTONOMO: nessun input o immagine esterna.
% Compilare con XeLaTeX due volte (anche su Overleaf).
% Font: Latin Modern Roman; arabo: Amiri se disponibile, altrimenti DejaVu Sans.
% Per togliere una lettura, cambiare {1} in {0} nel relativo blocco
% \\begin{lettura}{1}{Titolo} ... \\end{lettura}.
% La numerazione delle letture attive e l'indice si aggiornano automaticamente.
% Le citazioni rinviano alla bibliografia finale: nessun file .bib necessario.
\documentclass[11pt,a4paper]{article}
\usepackage[margin=18mm,top=20mm,bottom=20mm,headheight=15pt]{geometry}
\usepackage{fontspec}
\setmainfont{Latin Modern Roman}
\usepackage{polyglossia}
\setdefaultlanguage{italian}
\TeXXeTstate=1
\newcommand{\textarabic}[1]{\begingroup\arabicfont\beginR#1\endR\endgroup}
\IfFontExistsTF{Amiri}{\newfontfamily\arabicfont[Script=Arabic,Scale=1.05]{Amiri}}{\newfontfamily\arabicfont[Script=Arabic,Scale=.92]{DejaVu Sans}}
\usepackage{amsmath,amssymb}
\usepackage{microtype,multicol,enumitem,booktabs,tabularx,array}
\usepackage{xcolor,tikz,pgfplots,caption}
\usetikzlibrary{arrows.meta,positioning}
\pgfplotsset{compat=1.18}
\usepackage{titlesec,fancyhdr,needspace,xparse}
\usepackage[unicode,hidelinks]{hyperref}
\hypersetup{pdftitle={Economia: leggere e capire},pdfauthor={Materiale didattico di Marco Giovanni Ferrari}}
\definecolor{accent}{HTML}{25556B}
\setlength{\columnsep}{8mm}
\setlength{\parindent}{0pt}
\setlength{\parskip}{4.5pt plus 1pt}
\setlength{\emergencystretch}{2em}
\linespread{1.04}
\setlist[enumerate]{leftmargin=*,itemsep=4pt,topsep=4pt}
\titleformat{\section}{\Large\bfseries\color{accent}}{\thesection.}{.6em}{}
\titleformat{\subsection}{\normalsize\bfseries}{\thesubsection}{.5em}{}
\titlespacing*{\section}{0pt}{0pt}{8pt}
\titlespacing*{\subsection}{0pt}{11pt}{4pt}
\captionsetup{font=small,labelfont=bf,skip=5pt}
\pagestyle{fancy}\fancyhf{}
\fancyhead[L]{\small Economia: leggere e capire}
\fancyhead[R]{\small Letture e lessico}
\fancyfoot[C]{\thepage}
\renewcommand{\headrulewidth}{.3pt}
\newcommand{\fonte}[1]{\textsuperscript{[\ref{src:#1}]}}
\newcommand{\voce}[3]{\par\smallskip\noindent\begin{minipage}{\linewidth}\textbf{#1}\quad\textarabic{#3}\\[2pt]{\small #2.}\end{minipage}\par}
\NewDocumentEnvironment{lettura}{m m +b}{%
\ifnum#1=1\relax
\clearpage\section{#2}\begin{multicols}{2}\raggedcolumns
#3
\end{multicols}
\fi}{}
\newcommand{\attivita}{\small\setlength{\parskip}{3pt}\par\medskip\noindent\rule{\linewidth}{.4pt}\subsection*{Comprendere il testo}}
\newcommand{\lessico}{\subsection*{Parole da riutilizzare}}
\newcommand{\lingua}{\subsection*{Lavorare sulla lingua}}
\newcommand{\scrittura}{\subsection*{Rielaborare}}
\newcommand{\righe}{\par\vspace{3pt}\noindent\rule{\linewidth}{.2pt}\par\vspace{10pt}\noindent\rule{\linewidth}{.2pt}\par}
\begin{document}
\begin{titlepage}
\vspace*{22mm}
{\large Materiali di economia\par}
\vspace{12mm}
{\Huge\bfseries Economia:\\[5mm]leggere e capire\par}
\vspace{10mm}
{\Large Letture progressive per un anno scolastico\par}
\vspace{12mm}
\rule{\textwidth}{.5pt}
\vspace{5mm}
{\large Scuola secondaria superiore\\Italiano, lessico economico e comprensione del testo\par}
\vfill
{\large Materiali di Marco Giovanni Ferrari\par}
\vspace{4mm}
Versione di lavoro --- settembre 2026\\
Quattordici letture con esempi, grafici e attività.
\end{titlepage}
\section*{Come usare il libretto}
Questo percorso utilizza l'economia per imparare parole nuove, leggere testi sempre più lunghi e ricostruire un ragionamento. Le prime letture partono da situazioni quotidiane; le ultime collegano più argomenti e possono essere affrontate in diverse lezioni.

Leggi una prima volta per comprendere l'argomento generale. Poi torna ai paragrafi, cerca il significato delle parole nel contesto e utilizza il glossario. Infine rispondi alle domande e scrivi con parole tue. Non è necessario conoscere ogni parola prima di iniziare: molti significati diventano più chiari durante la lettura.

Le traduzioni arabe sono aiuti al significato usato nel testo, non sostituiscono tutte le possibili traduzioni di una parola. Dopo averle lette, prova a formulare la definizione in italiano. Le parole già incontrate ritornano nelle letture successive.

I numeri inventati per gli esercizi sono sempre presentati come esempi didattici. I pochi dati reali hanno una fonte e un periodo di riferimento. Le note numerate rinviano ai riferimenti finali. Le cifre storiche non sono quotazioni o aggiornamenti in tempo reale.

\subsection*{Tre abitudini di lettura}
\begin{enumerate}
\item Per ogni paragrafo, individua l'idea principale e una parola che la esprime.
\item Distingui un fatto, un esempio, un obiettivo e una previsione.
\item Quando leggi un numero, cerca l'unità, il periodo e il termine di confronto.
\end{enumerate}
\subsection*{Per il docente}
La progressione delle lunghezze riguarda il solo testo delle letture, esclusi titoli, attività e glossari. I testi più lunghi sono unità da distribuire su più incontri. Le consegne di scrittura possono essere svolte sul quaderno. Il sorgente permette di disattivare una lettura senza alterare manualmente l'indice. La suddivisione trimestrale è orientativa e non impone di completare una lettura in una singola lezione.

La raccolta riprende e rielabora i nuclei dei materiali del corso precedente: settori economici, reddito e patrimonio, beni pubblici, rischio e interesse, inflazione, crescita e trasformazioni economiche. Le altre letture sono raccordi e ampliamenti. Non è una trascrizione diplomatica dei fogli originali.
\clearpage
\tableofcontents
\clearpage\phantomsection\addcontentsline{toc}{part}{Primo trimestre --- Le parole fondamentali dell’economia}

% =================================================================
% INIZIO LETTURA 01 - circa 525 parole di testo

\begin{lettura}{1}{Bisogni, beni e servizi}


\subsection*{Una giornata di scelte}
Ogni mattina abbiamo bisogno di alcune cose. Beviamo acqua, mangiamo, indossiamo vestiti e ci spostiamo. Abbiamo anche bisogno di imparare, di parlare con altre persone e di riposare. Un \textbf{bisogno} è qualcosa che sentiamo necessario per la nostra vita o per il nostro benessere. Alcuni bisogni sono comuni a tutti. Altri cambiano con l'età, con le abitudini e con la situazione in cui viviamo.

Per soddisfare i bisogni usiamo delle \textbf{risorse}. Il tempo è una risorsa. Anche il denaro, l'acqua, la terra e il lavoro sono risorse. Le risorse disponibili non sono infinite. Una famiglia non può comprare tutto. Uno studente non può usare la stessa ora per fare due attività diverse. Per questo dobbiamo scegliere. L'economia studia anche come le persone usano risorse limitate per soddisfare bisogni diversi.

\subsection*{Beni e servizi}
Un \textbf{bene} materiale è un oggetto che possiamo utilizzare. Il pane, una penna e una bicicletta sono beni. Un \textbf{servizio} è un'attività svolta per soddisfare un bisogno. Una lezione, un taglio di capelli e un viaggio in autobus sono servizi. Quando compriamo il pane, riceviamo un bene. Quando paghiamo il biglietto dell'autobus, riceviamo un servizio di trasporto. In entrambi i casi soddisfiamo un bisogno.

Una stessa attività può offrire beni e servizi insieme. In un ristorante mangiamo del cibo, ma riceviamo anche un servizio: alcune persone cucinano, portano i piatti e puliscono il tavolo. Per capire la differenza dobbiamo osservare che cosa viene offerto. La distinzione è utile, anche quando nella vita quotidiana beni e servizi sono collegati.

\subsection*{Come usiamo i beni}
Un \textbf{bene di consumo} soddisfa direttamente un bisogno. Il pane che mangiamo è un bene di consumo. Un \textbf{bene strumentale} serve per produrre altri beni o servizi. Il forno di un panificio è un bene strumentale. Il panettiere usa il forno per preparare il pane che vende ai clienti.

La distinzione dipende dall'uso. Un computer usato per guardare un film a casa è un bene di consumo. Un computer usato in un ufficio per organizzare il lavoro è un bene strumentale. L'oggetto può essere simile, ma cambia la sua funzione. Per classificare un bene chiediamo quindi: chi lo usa e per fare che cosa?

\subsection*{Beni usati insieme o in alternativa}
Due beni sono \textbf{complementari} quando vengono usati insieme. Una stampante e le sue cartucce sono complementari. Anche una torcia e le batterie adatte sono complementari. Comprare una torcia senza avere le batterie può non essere sufficiente per ottenere la luce.

Due beni sono \textbf{sostituti} quando possiamo usare uno al posto dell'altro per soddisfare un bisogno simile. Due marche di quaderni possono essere sostitute. Per bere qualcosa durante la pausa possiamo scegliere tra due bibite diverse. Non è necessario comprare entrambe.

Immaginiamo infine che Mariam abbia denaro sufficiente per un quaderno oppure per una bibita. Mariam sceglie il quaderno perché deve scrivere durante la lezione. Rinuncia alla bibita, almeno per quel momento. La sua scelta mostra un'idea semplice: utilizzare una risorsa per uno scopo può impedire di utilizzarla per un altro. Nelle prossime letture incontreremo famiglie, imprese e Stati che devono compiere scelte di questo tipo.


Quando leggiamo una descrizione, cerchiamo prima il bisogno e poi il mezzo utilizzato per soddisfarlo. In questo modo possiamo riconoscere la stessa relazione anche quando cambiano gli oggetti dell'esempio.


\attivita
\begin{enumerate}

\item Quali risorse vengono nominate nella prima sezione?

\item Spiega la differenza tra il pane e un viaggio in autobus.

\item Perché il forno del panificio è un bene strumentale?

\item Quando due beni sono complementari? Trova l’esempio nel testo.

\end{enumerate}
\lessico

\voce{bisogno}{qualcosa sentito come necessario}{حاجة}

\voce{risorsa}{mezzo disponibile per fare qualcosa}{مورد}

\voce{bene}{oggetto utile a soddisfare un bisogno}{سلعة}

\voce{servizio}{attività svolta per soddisfare un bisogno}{خدمة}

\voce{consumo}{uso di beni e servizi per i propri bisogni}{استهلاك}

\voce{strumentale}{che serve per produrre}{إنتاجي}

\voce{complementari}{utilizzati insieme}{متكاملة}

\voce{sostituti}{utilizzabili uno al posto dell’altro}{بديلة}

\lingua
Completa con \emph{bene, servizio, risorsa}: il tempo è una \ldots; una lezione è un \ldots; una penna è un \ldots. Poi scrivi tre frasi con la struttura «\ldots serve per \ldots».
\scrittura
Scrivi 60--80 parole su una tua mattina. Nomina due beni, due servizi e una scelta che devi fare.

\end{lettura}
% FINE LETTURA 01

% =================================================================
% INIZIO LETTURA 02 - circa 614 parole di testo

\begin{lettura}{1}{La produzione e i settori economici}


\subsection*{Da una risorsa a un prodotto}
Il pane che troviamo in un negozio non compare da solo. Prima qualcuno coltiva il grano. Poi qualcuno trasforma il grano in farina. Un panettiere usa la farina per preparare il pane. Altre persone trasportano e vendono il prodotto. Tutte queste attività sono collegate. La \textbf{produzione} è l'insieme delle attività con cui otteniamo beni e servizi utili.

Per produrre servono risorse naturali, lavoro e strumenti. Il lavoro comprende attività manuali e attività intellettuali. Un agricoltore lavora la terra; un tecnico controlla una macchina; una persona organizza le consegne. Nessuna di queste attività basta sempre da sola. Le persone collaborano e svolgono compiti diversi.

\subsection*{Il settore primario}
Per descrivere le attività economiche le dividiamo in \textbf{settori}. Il settore primario comprende, nella classificazione scolastica tradizionale, attività che utilizzano direttamente le risorse naturali. L'agricoltura, l'allevamento, la pesca e la silvicoltura appartengono a questo settore. La silvicoltura riguarda la cura e l'utilizzo dei boschi. In questa classificazione anche l'estrazione di minerali viene spesso inserita nel primario.

I prodotti di queste attività possono essere consumati direttamente oppure diventare \textbf{materie prime}. Una materia prima viene utilizzata per ottenere un altro prodotto. Il cotone, per esempio, può diventare un tessuto. Il latte può essere bevuto oppure usato per produrre formaggio. L'uso successivo collega il settore primario agli altri settori.

\subsection*{Il settore secondario}
Il settore secondario comprende soprattutto la \textbf{trasformazione} delle materie prime e la costruzione di edifici e infrastrutture. Una fabbrica che produce tessuti e un'impresa che costruisce case operano in questo settore. Anche molte attività artigianali, come la produzione di mobili, appartengono al secondario.

Le macchine aiutano a trasformare i materiali, ma richiedono energia, manutenzione e persone preparate. In una fabbrica non lavorano soltanto gli operai che utilizzano le macchine. Servono anche tecnici, progettisti e persone che controllano la qualità. Il prodotto finale dipende dalla collaborazione tra questi lavoratori.

\subsection*{Il settore terziario}
Il settore terziario comprende i \textbf{servizi}. Trasporti, commercio, turismo, istruzione, sanità e attività bancarie sono esempi. Un autista offre un servizio anche se non produce un nuovo oggetto. Una scuola offre un servizio che permette agli studenti di imparare. I servizi hanno quindi valore economico e sociale.

Con l'espressione \textbf{terziario avanzato} indichiamo servizi che utilizzano molta conoscenza specialistica, come la ricerca e alcuni servizi informatici. Si tratta di una parte del terziario, non di un gruppo da sommare di nuovo al totale dei servizi. Inoltre, un'impresa tecnologica può svolgere attività differenti: produrre computer e progettare programmi non sono la stessa attività.

\subsection*{Una filiera e un paese}
La sequenza di attività necessarie per arrivare a un prodotto si chiama \textbf{filiera}. Nella filiera di una maglietta troviamo la coltivazione del cotone, la produzione del tessuto, la confezione, il trasporto e la vendita. Alcune attività possono avvenire in paesi diversi. Un oggetto semplice collega quindi molti lavoratori.

Anche in Egitto possiamo riconoscere attività agricole, industriali e di servizio: una coltivazione nella valle del Nilo, un laboratorio di mobili, un albergo o un servizio di trasporto. Questi esempi mostrano attività presenti nel paese; non indicano quale settore abbia il peso maggiore.

Per misurare il peso di un settore dobbiamo scegliere un indicatore. Possiamo contare i lavoratori oppure misurare il valore prodotto. Le due percentuali possono essere diverse. Inoltre, nelle statistiche internazionali l'estrazione mineraria è spesso compresa nell'industria. Prima di confrontare due grafici dobbiamo quindi leggere la classificazione e l'unità utilizzata. Una buona lettura dei dati comincia sempre dalle parole che li accompagnano.


Un'attività può inoltre cambiare nel tempo. Un negozio può iniziare a vendere anche attraverso internet e organizzare consegne. Il servizio resta collegato al commercio, ma richiede nuove competenze e nuovi strumenti. Le categorie aiutano a orientarsi; la descrizione concreta permette di capire quali lavori vengono svolti e come collaborano tra loro le persone. Per questo, negli esercizi, alla classificazione accompagneremo sempre una breve spiegazione.



\begin{center}\begin{minipage}{\linewidth}\centering
\begin{tikzpicture}[node distance=5mm, every node/.style={draw,rounded corners,align=center,text width=.82\linewidth,font=\small,inner sep=5pt}]
\node (a) {Coltivazione del cotone\\Settore primario};
\node[below=of a] (b) {Tessitura e confezione\\Settore secondario};
\node[below=of b] (c) {Trasporto e vendita\\Settore terziario};
\draw[-{Stealth}] (a)--(b);\draw[-{Stealth}] (b)--(c);
\end{tikzpicture}
\captionof{figure}{Una filiera semplificata: la maglietta.}
\end{minipage}\end{center}


\attivita
\begin{enumerate}

\item Ricostruisci il percorso del pane dalla coltivazione alla vendita.

\item Quali attività del secondario vengono nominate?

\item Perché il terziario avanzato non va contato due volte?

\item Perché un grafico sui lavoratori può essere diverso da un grafico sul valore prodotto?

\end{enumerate}
\lessico

\voce{produzione}{attività che permette di ottenere beni e servizi}{إنتاج}

\voce{settore}{gruppo di attività economiche}{قطاع}

\voce{materia prima}{materiale utilizzato per ottenere un prodotto}{مادة خام}

\voce{trasformazione}{passaggio da un materiale a un prodotto diverso}{تحويل}

\voce{filiera}{sequenza di attività collegate a un prodotto}{سلسلة الإنتاج}

\voce{manutenzione}{attività che mantiene un bene funzionante}{صيانة}

\voce{indicatore}{misura che descrive un fenomeno}{مؤشر}

\lingua
Trova nel testo i verbi collegati a queste parole: coltivazione, produzione, trasformazione, costruzione. Scrivi una frase con ciascun verbo.
\scrittura
Descrivi in 80--100 parole la filiera di un mobile. Usa «prima», «poi», «infine» e indica almeno un servizio.

\end{lettura}
% FINE LETTURA 02

% =================================================================
% INIZIO LETTURA 03 - circa 699 parole di testo

\begin{lettura}{1}{Reddito, patrimonio e risparmio}


\subsection*{Una fotografia e un periodo}
Quando diciamo che una persona ha molto denaro, possiamo riferirci a situazioni diverse. La persona può ricevere uno stipendio elevato ogni mese. Oppure può possedere una casa e molti risparmi, anche senza lavorare. Per descrivere queste situazioni usiamo due parole: \textbf{reddito} e \textbf{patrimonio}. Le parole sono collegate, ma non hanno lo stesso significato.

Il reddito è un flusso: si riferisce a un periodo. Possiamo parlare del reddito di un mese o di un anno. Il patrimonio, invece, si osserva in un momento preciso. È simile a una fotografia di ciò che una persona possiede. Dire «guadagna diecimila lire egiziane» non basta: dobbiamo sapere se le riceve ogni mese, ogni anno oppure una sola volta.

\subsection*{Da dove arriva il reddito}
Il \textbf{reddito da lavoro} deriva dall'attività svolta da una persona. Lo stipendio di un insegnante e il salario di un operaio sono esempi. Anche un professionista riceve un compenso per il proprio lavoro. Le somme possono essere diverse perché cambiano le ore lavorate, le competenze richieste, il tipo di contratto e molte altre condizioni.

Il \textbf{reddito da capitale} deriva dal possesso di risorse impiegate economicamente. Gli interessi su un deposito, i dividendi di un'impresa e gli affitti ricevuti sono esempi. Possedere una risorsa, però, non garantisce sempre un guadagno. Una casa può restare vuota e un'impresa può non distribuire dividendi. Il reddito da capitale può quindi variare.

Esistono anche entrate che non sono un pagamento per un lavoro svolto oggi o per un capitale investito. Una pensione o un sostegno pubblico sono esempi di \textbf{trasferimenti}. Nel bilancio di una famiglia queste somme sono entrate importanti. Per comprendere la loro origine è utile distinguerle dagli stipendi e dagli interessi. Un piccolo imprenditore, invece, può ricevere un reddito che combina il proprio lavoro e l'uso dei propri strumenti.

\subsection*{Che cosa comprende il patrimonio}
Nel patrimonio troviamo beni e diritti con valore economico: una casa, un terreno, un deposito bancario o una partecipazione in un'impresa. Per descrivere la situazione completa dobbiamo considerare anche i \textbf{debiti}, cioè somme che devono essere restituite. Il patrimonio netto è il valore delle attività possedute meno il valore dei debiti.

Immaginiamo una persona con una casa che vale un milione di lire egiziane, risparmi per centomila lire e un debito di quattrocentomila lire. Il valore delle attività è un milione e centomila lire. Sottraendo il debito otteniamo un patrimonio netto di settecentomila lire. Sono valori inventati per capire il calcolo, non prezzi di mercato.

Una casa ha valore, ma non può sempre essere trasformata rapidamente in denaro. Il denaro disponibile su un conto è generalmente più \textbf{liquido}. La liquidità riguarda la possibilità di usare una risorsa per pagare, oppure di convertirla in denaro senza grandi difficoltà e perdite. Una persona può avere un patrimonio elevato e poco denaro immediatamente disponibile.

\subsection*{Un mese con una spesa imprevista}
Riprendiamo la famiglia con duemila lire di risparmio mensile. Se deve pagare una riparazione straordinaria di tremila lire, quel mese le uscite aumentano. Può utilizzare una parte dei risparmi precedenti senza aver perso necessariamente il proprio reddito. La spesa imprevista modifica il bilancio del periodo. Per descriverla dobbiamo specificare se è occasionale oppure se si ripeterà ogni mese. Una difficoltà temporanea e una differenza stabile tra entrate e uscite richiedono infatti analisi diverse. Il bilancio aiuta a riconoscere questa distinzione prima di prendere nuove decisioni.

\subsection*{Consumo e risparmio}
Una famiglia usa le proprie entrate per pagare alimenti, trasporti, abitazione e altri consumi. La parte di reddito disponibile che non viene consumata costituisce il \textbf{risparmio}. Se le entrate disponibili sono dodicimila lire e i consumi sono diecimila, il risparmio è duemila lire. Il reddito disponibile è quello che resta dopo le imposte dirette e comprende i trasferimenti ricevuti.

Il risparmio può aumentare il patrimonio finanziario oppure servire a ridurre un debito. Tuttavia, il patrimonio cambia anche per altre ragioni: il valore di un immobile può salire o scendere. Perciò l'aumento del patrimonio non coincide sempre con il risparmio dell'anno.

Una famiglia può anche consumare più del reddito disponibile. In questo caso utilizza risparmi precedenti oppure si indebita. La scelta può essere temporanea, per esempio durante un periodo senza lavoro. Per capire la situazione non basta osservare un solo mese: bisogna conoscere le entrate, le spese, i debiti e la loro durata. Tenere un \textbf{bilancio} significa registrare queste informazioni e metterle in relazione.



\begin{center}\begin{minipage}{\linewidth}\centering\small
\begin{tabular}{lr}\toprule
Bilancio mensile immaginario & LE\\\midrule
Reddito disponibile & 12\,000\\
Consumi & 10\,000\\\midrule
Risparmio & 2\,000\\\bottomrule
\end{tabular}
\captionof{table}{Esempio didattico; LE indica lire egiziane.}
\end{minipage}\end{center}


\attivita
\begin{enumerate}

\item Perché il reddito ha bisogno di un periodo di riferimento?

\item Qual è la differenza tra stipendio, interesse e pensione?

\item Calcola il patrimonio netto dell’esempio.

\item Come può una famiglia avere un patrimonio elevato e poca liquidità?

\item Se una famiglia risparmia, il valore totale del suo patrimonio aumenta necessariamente della stessa somma?

\end{enumerate}
\lessico

\voce{reddito}{flusso di risorse ricevute in un periodo}{دخل}

\voce{patrimonio}{insieme di beni e diritti con valore economico}{ثروة}

\voce{debito}{somma da restituire}{دَين}

\voce{risparmio}{reddito disponibile non consumato}{ادخار}

\voce{trasferimento}{entrata senza una produzione corrente in cambio}{تحويل}

\voce{liquidità}{disponibilità di mezzi utilizzabili per pagare}{سيولة}

\voce{bilancio familiare}{registrazione delle entrate e delle uscite}{ميزانية الأسرة}

\lingua
Completa quattro frasi usando «mentre»: «Il reddito si misura in un periodo, mentre…». Poi trova nel testo tre coppie di parole di significato opposto.
\scrittura
Scrivi 100 parole per descrivere un bilancio familiare immaginario. Dichiara il periodo e distingui entrate, consumi e risparmio.

\end{lettura}
% FINE LETTURA 03

% =================================================================
% INIZIO LETTURA 04 - circa 829 parole di testo

\begin{lettura}{1}{La ricchezza di un paese: il PIL}


\subsection*{Dalla famiglia all'economia nazionale}
Nella lettura precedente abbiamo osservato il reddito di una famiglia. Ora allarghiamo lo sguardo a un paese. In un anno, milioni di persone producono beni e servizi. Alcune coltivano, altre costruiscono, insegnano o trasportano merci. Per descrivere questa attività economica utilizziamo il \textbf{prodotto interno lordo}, chiamato PIL. In inglese troviamo la sigla GDP, che significa \emph{Gross Domestic Product}.

Il PIL misura il valore dei beni e servizi finali prodotti all'interno di un paese durante un periodo, generalmente un anno o un trimestre. Ogni parola è importante. «Prodotto» indica che osserviamo un'attività di produzione. «Interno» indica il territorio economico considerato. «Lordo» significa che non abbiamo sottratto il consumo del capitale fisso: macchinari, edifici e altri strumenti si usurano durante la produzione.\fonte{pil}

Il PIL non è il patrimonio del paese e non è il denaro posseduto dal governo. È una misura della produzione in un periodo. Se una casa costruita molti anni fa viene venduta oggi, il suo intero prezzo non rappresenta una nuova produzione di quest'anno. Possono invece entrare nel PIL i servizi prodotti oggi per realizzare la vendita.

\subsection*{Perché non sommiamo tutto}
Consideriamo una filiera semplificata. Un agricoltore vende grano per venti lire a un mulino. Il mulino vende farina per trenta lire a un panificio. Il panificio vende il pane ai consumatori per cinquanta lire. Se sommiamo venti, trenta e cinquanta, otteniamo cento. Ma stiamo contando più volte il valore del grano e della farina contenuti nel pane.

In questo esempio, senza altri acquisti intermedi e senza imposte sui prodotti, il valore del bene finale è cinquanta. Possiamo ottenere lo stesso risultato sommando il \textbf{valore aggiunto} delle tre attività. L'agricoltore aggiunge venti, il mulino aggiunge dieci e il panificio aggiunge venti. Il valore aggiunto è la differenza tra il valore della produzione e quello dei beni e servizi intermedi utilizzati.

Questa differenza non coincide con il profitto. Il valore aggiunto deve anche remunerare il lavoro e comprende altre componenti. Gli stipendi, per esempio, non sono acquisti intermedi come la farina. Distinguere queste parole permette di capire perché un'impresa può creare molto valore e avere comunque costi elevati.

\subsection*{Il PIL per abitante}
Un paese molto popoloso può produrre più di un paese piccolo. Per confrontare la produzione media per persona calcoliamo il \textbf{PIL pro capite}: dividiamo il PIL totale per il numero di abitanti. L'espressione latina «pro capite» significa «per persona». Il risultato è una media, non la somma ricevuta effettivamente da ogni abitante.

Immaginiamo due paesi. Il paese A produce cento miliardi di unità monetarie e ha dieci milioni di abitanti. Il suo PIL pro capite è diecimila. Il paese B produce sessanta miliardi e ha tre milioni di abitanti: il PIL pro capite è ventimila. A produce di più in totale; B produce di più per abitante. Sono due risposte diverse a due domande diverse.

Anche all'interno di un paese la media può nascondere grandi differenze. Se due persone ricevono rispettivamente due e diciotto unità, la media è dieci. Nessuna delle due riceve dieci. Per conoscere la distribuzione dobbiamo aggiungere altre informazioni. Il PIL pro capite, da solo, non descrive le disuguaglianze.

\subsection*{Quantità, prezzi e confronti}
Il valore della produzione può aumentare perché produciamo più beni oppure perché aumentano i prezzi. Se un laboratorio produce cento sedie a cento lire, il valore è diecimila. Se l'anno successivo produce ancora cento sedie a centoventi lire, il valore sale a dodicimila, ma il numero di sedie non cambia.

Il \textbf{PIL nominale} utilizza i prezzi del periodo osservato. Il \textbf{PIL reale} cerca di separare la variazione delle quantità da quella dei prezzi. Quando leggiamo che un'economia cresce, dobbiamo verificare quale misura viene usata. Nelle analisi della crescita si osserva generalmente una misura reale.

Anche il confronto tra paesi richiede attenzione. I paesi possono usare monete differenti e avere prezzi differenti. Convertire tutto in dollari è utile, ma il tasso di cambio può cambiare molto. Esistono anche confronti che tengono conto delle differenze nel costo di un insieme di beni e servizi: sono confronti a parità di potere d'acquisto.

\subsection*{Il periodo e il confine}
Un servizio prodotto in un paese entra nella sua produzione interna anche se l'impresa appartiene a proprietari esteri, secondo le regole della contabilità nazionale. La nazionalità del proprietario e il luogo della produzione sono informazioni diverse. Anche il periodo conta: confrontare un trimestre con un anno intero senza adattare le misure può portare a un errore. Prima di leggere il valore, controlliamo sempre il titolo della tabella e le note.

\subsection*{Produzione e benessere}
Il PIL fornisce informazioni importanti, ma non misura ogni aspetto della vita. Non descrive direttamente la salute, la qualità dell'aria, il tempo libero o la sicurezza. Inoltre, molte attività domestiche non pagate, come prendersi cura di un familiare, non sono comprese nello stesso modo dei servizi venduti sul mercato.

Per valutare il \textbf{benessere} affianchiamo quindi al PIL altri indicatori. Possiamo osservare l'accesso alla scuola, le condizioni delle abitazioni e la distribuzione dei redditi. Non dobbiamo scegliere una sola misura per rispondere a tutte le domande. Dobbiamo scegliere la misura adatta alla domanda che stiamo facendo e leggere con attenzione ciò che quella misura include.



\begin{center}\begin{minipage}{\linewidth}\centering\small
\begin{tabular}{lrr}\toprule
 & Paese A & Paese B\\\midrule
PIL (miliardi) & 100 & 60\\
Abitanti (milioni) & 10 & 3\\
PIL pro capite & 10\,000 & 20\,000\\\bottomrule
\end{tabular}
\captionof{table}{Due paesi immaginari, stessa unità monetaria.}
\end{minipage}\end{center}


\attivita
\begin{enumerate}

\item Che cosa misura il PIL e che cosa non misura?

\item Perché nell’esempio del pane non dobbiamo sommare 20, 30 e 50?

\item Quale paese immaginario ha il PIL totale maggiore? Quale ha il PIL pro capite maggiore?

\item Nell’esempio delle sedie aumenta la quantità prodotta?

\item Indica due aspetti del benessere che richiedono informazioni ulteriori.

\end{enumerate}
\lessico

\voce{PIL}{valore della produzione finale interna in un periodo}{الناتج المحلي الإجمالي}

\voce{valore aggiunto}{valore prodotto meno acquisti intermedi}{قيمة مضافة}

\voce{pro capite}{per abitante}{للفرد}

\voce{media}{totale diviso per il numero di osservazioni}{متوسط}

\voce{nominale}{misurato ai prezzi del periodo}{اسمي}

\voce{reale}{corretto per la variazione dei prezzi}{حقيقي}

\voce{benessere}{condizione di vita delle persone}{رفاه}

\lingua
Riscrivi con parole tue «una media non descrive la distribuzione». Poi completa: «Il valore aumenta perché…, ma la quantità…».
\scrittura
In 100--120 parole spiega a un compagno perché «il paese con il PIL più alto è quello dove tutti vivono meglio» non è una conclusione sufficiente.

\end{lettura}
% FINE LETTURA 04

% =================================================================
% INIZIO LETTURA 05 - circa 935 parole di testo

\begin{lettura}{1}{Beni pubblici e ruolo dello Stato}


\subsection*{Un servizio per molte persone}
Pensiamo a una strada illuminata. La luce permette a molte persone di camminare con maggiore sicurezza. Chi passa per primo non porta via tutta la luce: anche chi arriva dopo può utilizzarla. Inoltre, è difficile illuminare soltanto il percorso delle persone che hanno pagato. Queste caratteristiche ci aiutano a comprendere il significato economico di \textbf{bene pubblico}.

Un bene pubblico puro presenta due proprietà: non rivalità e non escludibilità. La \textbf{non rivalità} significa che l'utilizzo da parte di una persona non riduce la possibilità di utilizzo da parte delle altre. La \textbf{non escludibilità} significa che impedire l'accesso a chi non paga è impossibile o molto difficile e costoso. Le due caratteristiche devono essere considerate insieme.

\subsection*{Pubblico non significa soltanto statale}
Un panino è rivale: se lo mangio io, un'altra persona non può mangiare lo stesso panino. È anche escludibile: il venditore può consegnarlo soltanto a chi paga. Una trasmissione a pagamento può invece essere non rivale, entro certi limiti tecnici, ma escludibile attraverso un abbonamento. Questi esempi mostrano che rivalità ed escludibilità sono proprietà differenti.

Un servizio offerto dallo Stato non è necessariamente un bene pubblico puro. Una visita medica richiede il tempo di un medico. Se il medico visita una persona, nello stesso momento non può visitare molte altre persone. Anche i posti in una classe sono limitati. Sanità e istruzione possono essere finanziate pubblicamente per molte ragioni, ma non diventano per questo beni perfettamente non rivali.

Le caratteristiche possono inoltre cambiare con le condizioni. Una strada poco frequentata permette a molti veicoli di passare senza ostacolarsi. Se il traffico aumenta molto, ogni automobile rallenta le altre. Compare la \textbf{congestione}. Per classificare correttamente un bene dobbiamo osservare la situazione concreta, non soltanto il suo nome.

\subsection*{Chi paga per qualcosa che tutti possono usare?}
Immaginiamo un quartiere che voglia finanziare l'illuminazione di una piazza. Ogni famiglia beneficia della luce, ma una famiglia può pensare: «Non contribuisco; forse pagheranno gli altri». Se molte famiglie ragionano così, i contributi non bastano. La piazza rimane buia anche se molte persone vorrebbero averla illuminata.

Questo comportamento viene chiamato problema del \textbf{free rider}, cioè di chi beneficia di qualcosa senza contribuire al suo finanziamento. Non significa che ogni persona si comporti sempre in questo modo. Significa che la non escludibilità può rendere difficile raccogliere contributi volontari sufficienti. Le associazioni e la collaborazione possono aiutare, ma non eliminano automaticamente il problema.

Quando le risorse destinate a un'attività sono inferiori a quelle utili per la collettività, parliamo di \textbf{sottoinvestimento}. Un'impresa privata può avere difficoltà a vendere un bene se non può escludere chi non paga. Per questo lo Stato e gli enti locali possono organizzare il finanziamento e la produzione di alcuni beni pubblici. Il finanziamento pubblico può anche accompagnarsi alla produzione affidata a imprese esterne.

\subsection*{Entrate pubbliche e spesa pubblica}
Per finanziare servizi e attività lo Stato raccoglie entrate. Una parte importante deriva dalle \textbf{imposte}, pagamenti obbligatori stabiliti dalle norme. Nel linguaggio quotidiano la parola «tasse» viene spesso usata per tutti questi pagamenti. Nel lessico tributario italiano, invece, una tassa è collegata a uno specifico servizio o atto pubblico, mentre un'imposta non è il prezzo diretto di un servizio individuale. Le denominazioni giuridiche possono variare tra paesi.

La \textbf{spesa pubblica} comprende impieghi diversi delle risorse: stipendi, acquisti, infrastrutture e trasferimenti. Una strada richiede materiali e lavoro oggi, ma può offrire servizi per molti anni. Un sostegno a una famiglia è invece un trasferimento: permette alla famiglia di disporre di risorse senza vendere direttamente un nuovo bene allo Stato.

Le risorse pubbliche sono limitate. Finanziare un progetto può ridurre le risorse disponibili per un altro. Per questo un bilancio pubblico deve considerare costi, obiettivi e persone coinvolte. Non basta dire che un servizio è utile. Bisogna chiedere quanto costa, come viene finanziato e se raggiunge le persone per cui è stato previsto.

\subsection*{Effetti su persone esterne allo scambio}
Un'officina e un cliente possono accordarsi sul prezzo di una riparazione. Se l'officina produce molto rumore, anche i vicini subiscono una conseguenza. I vicini, però, non partecipano allo scambio. Un effetto su persone esterne alla transazione, non pienamente considerato nel prezzo, si chiama \textbf{esternalità}.

L'esternalità può essere negativa, come il rumore, oppure positiva. La cura di un giardino, per esempio, può rendere più piacevole anche la strada vicina. La parola «esterna» non significa che l'effetto avviene fuori dal paese. Significa che coinvolge qualcuno esterno alla decisione o allo scambio considerato.

Pensiamo al traffico in una grande città come il Cairo. Un viaggio offre un beneficio a chi si sposta. I gas di scarico e il rumore possono però influenzare anche altre persone. Per discutere il problema dobbiamo distinguere il beneficio individuale dai costi collettivi. Torneremo su questa distinzione nella lettura sull'ambiente.

\subsection*{Dalla domanda al progetto}
Torniamo alla piazza iniziale. Le famiglie possono discutere non soltanto se desiderano la luce, ma anche quante lampade servono e quali orari sono utili. Un progetto richiede una scelta tecnica e una scelta di finanziamento. Dopo l'installazione occorre pagare energia e manutenzione. Il costo iniziale non esaurisce quindi l'impegno. Se una lampada si rompe, qualcuno deve essere responsabile della riparazione. Questo esempio mostra che garantire un servizio significa organizzare un'attività nel tempo: raccogliere risorse, utilizzare strumenti, controllare il funzionamento e intervenire quando serve.

\subsection*{Osservare gli strumenti}
Le autorità possono utilizzare regole, controlli, imposte, incentivi o servizi alternativi. Una regola può limitare un'attività inquinante; un trasporto pubblico affidabile può offrire un'alternativa. Strumenti diversi agiscono in modi diversi. La loro efficacia dipende anche dalla possibilità di applicarli e dalle condizioni delle persone.

Studiare il ruolo dello Stato significa quindi riconoscere problemi, strumenti e risultati. Dobbiamo distinguere l'obiettivo di una politica dagli effetti che produce realmente. Possiamo comprendere perché un intervento viene proposto e, allo stesso tempo, chiedere quali informazioni servano per valutarlo. Questa distinzione sarà utile anche quando leggeremo delle riforme economiche egiziane.


\attivita
\begin{enumerate}

\item Quali due proprietà definiscono un bene pubblico puro?

\item Perché una visita medica non è un bene pubblico puro?

\item Ricostruisci il problema del quartiere che deve finanziare la luce.

\item Che cosa significa esternalità?

\item Distingui un obiettivo pubblico da uno strumento per raggiungerlo.

\end{enumerate}
\lessico

\voce{rivalità}{l’uso di uno limita l’uso degli altri}{تزاحم في الاستهلاك}

\voce{escludibilità}{possibilità di impedire l’accesso}{إمكانية الاستبعاد}

\voce{congestione}{eccesso di utilizzatori che si ostacolano}{ازدحام}

\voce{imposta}{pagamento obbligatorio non legato a un singolo servizio}{ضريبة}

\voce{spesa pubblica}{risorse spese dalle amministrazioni pubbliche}{إنفاق عام}

\voce{esternalità}{effetto su soggetti esterni allo scambio}{أثر خارجي}

\voce{incentivo}{strumento che incoraggia un comportamento}{حافز}

\lingua
Trova le frasi costruite con «se». Riscrivine due cambiando l’esempio ma conservando la relazione di causa e conseguenza.
\scrittura
Scrivi 120--150 parole su un servizio del tuo quartiere: chi lo usa, chi lo finanzia, quali risorse richiede? Evita di chiamarlo bene pubblico senza verificare le due proprietà.

\end{lettura}
% FINE LETTURA 05

\clearpage\phantomsection\addcontentsline{toc}{part}{Secondo trimestre --- Moneta e scelte economiche}

% =================================================================
% INIZIO LETTURA 06 - circa 510 parole di testo

\begin{lettura}{1}{La moneta, le banche e il credito}


\subsection*{La moneta facilita gli scambi}
Immaginiamo un falegname che voglia comprare del pane. Senza moneta dovrebbe offrire in cambio una sedia e trovare un panettiere interessato proprio a quella sedia. Lo scambio sarebbe difficile: le due persone dovrebbero desiderare, nello stesso momento, ciò che l'altra offre.

La moneta evita questo problema perché è accettata da molte persone come \textbf{mezzo di scambio}. Il falegname può vendere la sedia a un cliente, ricevere moneta e usare una parte della somma per comprare il pane. Il panettiere accetta il pagamento perché sa che potrà utilizzare quella moneta per altri acquisti. La moneta collega quindi scambi diversi, tra persone diverse e in momenti diversi.

La moneta può essere anche una \textbf{riserva di valore}. Chi la riceve non deve spenderla immediatamente: può conservarla per un acquisto futuro. Questa funzione non è perfetta. Se nel tempo i prezzi aumentano, la stessa somma permette di comprare meno beni. Approfondiremo questo problema nella lettura sull'inflazione.

\subsection*{Risparmiatori e investitori}
Non tutte le persone utilizzano il proprio denaro nello stesso momento. Un \textbf{risparmiatore} rinuncia a spendere oggi una parte del proprio reddito e la conserva per il futuro. Può anche decidere di mettere questa somma a disposizione di altri, con l'obiettivo di ricevere un interesse. In questo caso accetta un rischio: il denaro potrebbe non essere restituito completamente o nei tempi previsti.

Un \textbf{investitore}, per esempio un'impresa, può invece avere bisogno di denaro oggi. Vuole acquistare una macchina, aprire un negozio o aumentare la produzione, ma non possiede ancora tutte le risorse necessarie. Può chiedere un prestito e accettare di restituire in futuro la somma ricevuta, aggiungendo un interesse. L'interesse rappresenta quindi un guadagno per chi mette a disposizione il risparmio e un costo per chi riceve il finanziamento.

\subsection*{La banca come intermediario}
Risparmiatori e investitori spesso non si conoscono. Inoltre, un singolo risparmiatore può avere una somma piccola, mentre un'impresa può aver bisogno di un finanziamento più grande. La banca svolge una funzione di \textbf{intermediazione}: raccoglie risorse dai risparmiatori e concede prestiti a famiglie e imprese.

La banca deve valutare il rischio che un prestito non venga restituito. Anche il risparmiatore, quando sceglie come utilizzare il proprio denaro, deve considerare il rapporto tra possibile guadagno e rischio. L'intermediazione non elimina l'incertezza, ma permette di organizzare il passaggio di risorse da chi vuole conservarle e farle fruttare a chi vuole utilizzarle per un progetto.

Pensiamo a un panificio che abbia bisogno di 100.000 LE per acquistare un nuovo forno. Molti risparmiatori depositano in banca somme più piccole. La banca utilizza le risorse disponibili per finanziare il panificio, dopo aver esaminato la sua richiesta. Il panificio paga un interesse sul prestito; la banca può riconoscere un interesse ai risparmiatori. Una parte della differenza serve alla banca per sostenere i propri costi e affrontare i rischi dell'attività.

Il percorso può essere riassunto così: alcune persone risparmiano, la banca raccoglie e organizza le risorse, altre persone le usano per investire. Il denaro si sposta verso un progetto che potrebbe produrre nuovi beni, servizi e redditi. Il risultato, però, non è garantito: proprio per questo il rischio e l'interesse saranno il centro della prossima lettura.





\attivita
\begin{enumerate}

\item Perché la moneta facilita lo scambio tra il falegname e il panettiere?

\item Che cosa significa riserva di valore?

\item Quali esigenze differenti hanno il risparmiatore e l’investitore?

\item Quale funzione svolge la banca tra questi due soggetti?

\item Chi paga l’interesse e chi può riceverlo?

\item Perché l’intermediazione non elimina il rischio?

\end{enumerate}
\lessico

\voce{moneta}{mezzo accettato per pagare e conservare valore}{نقود}

\voce{risparmiatore}{chi conserva una parte del proprio reddito}{مدّخر}

\voce{investitore}{chi impiega risorse in un progetto}{مستثمر}

\voce{intermediazione}{collegamento tra soggetti con esigenze diverse}{وساطة}

\voce{prestito}{somma ricevuta che deve essere restituita}{قرض}

\voce{interesse}{somma pagata o ricevuta per un finanziamento}{فائدة}

\voce{rischio}{possibilità che il risultato sia diverso da quello previsto}{مخاطرة}

\lingua
Completa quattro frasi usando correttamente «oggi», «in futuro», «perché» e «quindi». Poi costruisci la famiglia di parole: risparmiare, risparmio, risparmiatore.
\scrittura
Riassumi il caso del panificio in 100--130 parole. Usa le parole risparmiatore, banca, investitore, prestito, interesse e rischio.

\end{lettura}
% FINE LETTURA 06

% =================================================================
% INIZIO LETTURA 07 - circa 545 parole di testo

\begin{lettura}{1}{Interesse, rischio e rendimento}


\subsection*{Un titolo annuale senza cedola}
Un'impresa ha bisogno di denaro e decide di emettere un \textbf{titolo} della durata di un anno. Un risparmiatore acquista il titolo pagando 1.000 LE. In questo modo presta 1.000 LE all'impresa. Dopo un anno l'impresa si impegna a restituire 1.050 LE.

Il titolo è \textbf{senza cedola}: durante l'anno non sono previsti pagamenti intermedi. L'intera somma viene pagata alla scadenza. Le 1.000 LE versate all'inizio sono il \textbf{capitale iniziale}. Le 50 LE aggiuntive sono l'\textbf{interesse}. La somma finale di 1.050 LE è il \textbf{montante}.

Il \textbf{tasso di interesse} misura l'interesse in percentuale rispetto al capitale iniziale. Nell'esempio:
\[
\begin{aligned}
I&=C\cdot i=1.000\cdot 0{,}05=50\text{ LE},\\
M&=C+I=1.050\text{ LE}.
\end{aligned}
\]
Il capitale è indicato con $C$, l'interesse con $I$, il tasso con $i$ e il montante con $M$. Dire che il tasso è del 5\% annuo significa che, per ogni 100 LE di capitale, sono previste 5 LE di interesse dopo un anno.

Facciamo un secondo esempio. Un titolo annuale ha un capitale iniziale di 2.000 LE e un tasso del 4\%. L'interesse è $2.000\cdot 0{,}04=80$ LE; il montante è quindi 2.080 LE. Se invece conosciamo un capitale di 5.000 LE e un montante di 5.300 LE, l'interesse è 300 LE e il tasso è $300/5.000=0{,}06$, cioè 6\%.

\subsection*{Rischio e interesse}
Il pagamento finale è previsto dal titolo, ma non è sempre certo. L'impresa potrebbe avere difficoltà e non riuscire a restituire tutto il denaro. Questa possibilità rappresenta un \textbf{rischio} per chi acquista il titolo.

In generale, un investitore accetta un rischio maggiore soltanto se può ottenere un interesse maggiore. Per esempio, tra due titoli annuali simili, uno considerato più sicuro potrebbe offrire il 4\% e uno considerato più rischioso il 7\%. Il maggiore interesse è una compensazione per il maggiore rischio. Non è però una garanzia di guadagno: proprio il titolo che promette di più potrebbe non essere rimborsato completamente.

\subsection*{Lo spread: misurare una differenza}
Lo \textbf{spread} è la differenza tra due tassi o rendimenti confrontabili. Se il titolo A offre il 7\% e il titolo B offre il 4\%, lo spread è $7\%-4\%=3$ punti percentuali. Si può anche dire che lo spread è di 300 punti base, perché un punto percentuale corrisponde a 100 punti base.

Un altro esempio: tra un tasso dell'8,5\% e uno del 6\% lo spread è di 2,5 punti percentuali, cioè 250 punti base. Uno spread maggiore può indicare che il mercato percepisce un rischio maggiore, ma misura soltanto una differenza: non spiega da solo tutte le sue cause.

\subsection*{Titolo e azione: due posizioni diverse}
Chi compra un titolo presta denaro ed è un creditore. Chi compra un'\textbf{azione}, invece, diventa proprietario di una piccola parte dell'impresa e partecipa ai suoi risultati. L'azionista può ricevere dividendi, ma questi dipendono dalle decisioni e dai risultati dell'impresa; non sono un interesse stabilito come nel nostro esempio.

La differenza è importante anche quando un'impresa non riesce a pagare tutti. In generale, i creditori hanno la priorità nel rimborso rispetto agli azionisti. Prima vengono soddisfatti i crediti secondo l'ordine stabilito dalle regole; gli azionisti ricevono soltanto ciò che eventualmente rimane. Per questo un'azione può offrire un guadagno elevato quando l'impresa va bene, ma sopporta per ultima le perdite quando le risorse non bastano.





\attivita
\begin{enumerate}

\item Nel primo esempio, quali sono capitale iniziale, interesse, tasso e montante?

\item Calcola interesse e montante di 3.000 LE al 5\% per un anno.

\item Perché un titolo con interesse maggiore può essere più rischioso?

\item Calcola lo spread tra 9\% e 6,5\% in punti percentuali e punti base.

\item Qual è la differenza principale tra un titolo e un’azione?

\item Chi ha generalmente priorità nel rimborso: il creditore o l’azionista?

\end{enumerate}
\lessico

\voce{titolo}{documento finanziario che rappresenta un rapporto economico}{ورقة مالية}

\voce{capitale iniziale}{somma impiegata all’inizio}{رأس المال الابتدائي}

\voce{interesse}{somma aggiuntiva pagata sul capitale}{فائدة}

\voce{tasso di interesse}{interesse espresso in percentuale del capitale}{سعر الفائدة}

\voce{montante}{capitale iniziale più interesse}{المبلغ المتراكم}

\voce{rischio}{possibilità di non ottenere il risultato previsto}{مخاطرة}

\voce{spread}{differenza tra due tassi o rendimenti}{فارق}

\voce{azione}{quota di proprietà di un’impresa}{سهم}

\voce{priorità}{diritto di essere soddisfatto prima di altri}{أولوية}

\lingua
Trova nel testo le espressioni «dopo un anno», «durante l’anno», «alla scadenza» e «all’inizio». Usale per descrivere in ordine il primo esempio.
\scrittura
Spiega in 100--130 parole il primo titolo e il rapporto tra rischio e interesse. Concludi distinguendo il creditore dall’azionista.

\end{lettura}
% FINE LETTURA 07

% =================================================================
% INIZIO LETTURA 08 - circa 1150 parole di testo

\begin{lettura}{1}{Le disuguaglianze economiche}

\subsection*{La distribuzione: media e mediana}
In una società, la ricchezza e il reddito non sono distribuiti in modo uguale. Alcune persone guadagnano molto, altre poco; alcune possiedono case, terreni o imprese, mentre altre non hanno un patrimonio significativo. Queste differenze vengono chiamate \textbf{disuguaglianze economiche}. Osservarle è importante, ma non basta sapere che esistono: bisogna anche capire come si misurano, da dove derivano e quanto sia possibile cambiare la propria posizione economica nel corso della vita.

La \textbf{distribuzione del reddito} descrive come il reddito complessivo di una società è diviso tra le persone o le famiglie. Per analizzarla si utilizzano spesso due valori: la \textbf{media} e la \textbf{mediana}.

La media si calcola sommando tutti i redditi e dividendo il risultato per il numero delle persone. Consideriamo questi cinque redditi mensili:
\[
10,\quad 10,\quad 10,\quad 10,\quad 10.
\]
La somma è 50. Dividendo per cinque persone otteniamo un reddito medio di 10.

Consideriamo ora un secondo gruppo:
\[
5,\quad 5,\quad 5,\quad 5,\quad 30.
\]
Anche in questo caso la somma è 50 e la media è 10. La situazione economica dei due gruppi, però, è molto diversa. Nel primo gruppo tutti ricevono la stessa somma. Nel secondo, una persona riceve molto più delle altre.

Per comprendere meglio questa differenza possiamo utilizzare la mediana. La mediana è il valore che si trova al centro dopo aver ordinato i dati dal più piccolo al più grande. Nel primo gruppo la mediana è 10; nel secondo è 5. La media può quindi essere influenzata da pochi redditi molto alti. La mediana descrive meglio il reddito della persona che si trova al centro della distribuzione. Per questo, quando si confrontano le condizioni economiche di due paesi, è utile osservare entrambi i valori.

\begin{center}
\begin{minipage}{\linewidth}\centering\small
\begin{tabular}{lcc}
\toprule
 & Gruppo A & Gruppo B \\
\midrule
Reddito totale & 50 & 50 \\
Reddito medio & 10 & 10 \\
Reddito mediano & 10 & 5 \\
\bottomrule
\end{tabular}
\captionof{table}{La stessa media può descrivere distribuzioni differenti.}
\end{minipage}
\end{center}

\subsection*{Redditi da lavoro e redditi da capitale}
Le persone possono ottenere un reddito da fonti differenti. Una prima distinzione importante è quella tra \textbf{reddito da lavoro} e \textbf{reddito da capitale}.

Il reddito da lavoro è il compenso ricevuto per un'attività lavorativa. Ne sono esempi lo stipendio di un insegnante, il salario di un operaio e il compenso di un tecnico. Le differenze tra i redditi da lavoro possono dipendere da numerosi fattori: il numero di ore lavorate, il livello di formazione, l'esperienza, la responsabilità richiesta e la domanda di una determinata professione. Possono inoltre dipendere dal settore economico, dalla produttività delle imprese, dal territorio, dai contratti e dalla capacità dei lavoratori di negoziare il proprio salario.

Un reddito più alto non è quindi necessariamente il risultato di un maggiore impegno personale. Le condizioni di partenza e le caratteristiche del mercato del lavoro possono avere un ruolo altrettanto importante. In alcuni casi possono intervenire anche forme di discriminazione legate, per esempio, al genere, all'origine o alla condizione sociale.

Il reddito da capitale deriva invece dal possesso di una risorsa economica. Una persona può ricevere un interesse perché ha prestato del denaro, un affitto perché possiede un'abitazione oppure un dividendo perché possiede una parte di un'impresa.

Anche l'aumento del valore di un bene può produrre un guadagno. Per esempio, una persona può acquistare un immobile o un'azione e rivenderli successivamente a un prezzo più alto. Tuttavia, il valore può anche diminuire: possedere capitale permette di ottenere un rendimento, ma comporta generalmente anche un rischio.

La differenza principale è quindi questa: il reddito da lavoro richiede un'attività lavorativa, mentre il reddito da capitale richiede il possesso di una risorsa. Una persona può ricevere entrambe le forme di reddito. Un lavoratore può, per esempio, ricevere uno stipendio e utilizzare una parte dei propri risparmi per acquistare titoli o azioni.

\subsection*{Perché le disuguaglianze possono aumentare}
I redditi da capitale hanno un ruolo importante nella formazione delle disuguaglianze. Chi possiede un patrimonio può investirlo e ottenere nuovi redditi. Una parte di questi redditi può essere nuovamente investita, facendo crescere ulteriormente il patrimonio.

Chi non possiede risorse iniziali ha meno possibilità di iniziare questo processo. In alcuni casi deve utilizzare quasi tutto il proprio reddito per soddisfare i bisogni quotidiani e non riesce a risparmiare. Anche l'eredità può aumentare le differenze. Alcune persone ricevono dalla propria famiglia una casa, un'impresa o una somma di denaro. Altre iniziano la vita adulta senza un patrimonio o con debiti da restituire.

Le differenze di reddito e di patrimonio possono quindi dipendere da molti elementi:
\begin{itemize}
\item il tipo di lavoro e il livello del salario;
\item la proprietà di case, terreni, imprese o strumenti finanziari;
\item la possibilità di risparmiare e investire;
\item l'eredità ricevuta dalla famiglia;
\item l'accesso all'istruzione e alle cure sanitarie;
\item il luogo nel quale una persona nasce e vive;
\item eventi imprevisti, come una malattia o la perdita del lavoro.
\end{itemize}

La disuguaglianza economica non nasce dunque da una sola causa. È il risultato dell'incontro tra scelte personali, condizioni familiari, funzionamento dei mercati e istituzioni pubbliche.

\subsection*{Mobilità e uguaglianza delle opportunità}
Per valutare le opportunità offerte da una società non è sufficiente osservare la distribuzione del reddito in un singolo momento. È utile anche chiedersi se le persone possano cambiare la propria posizione economica.

La \textbf{mobilità economica} indica la possibilità di passare da una condizione economica a un'altra. Si parla di mobilità verso l'alto quando una persona migliora la propria posizione e di mobilità verso il basso quando la peggiora. La mobilità può essere osservata all'interno della vita di una persona. Un giovane può iniziare con un lavoro poco pagato e ottenere in seguito una posizione migliore. Al contrario, un lavoratore può perdere il proprio impiego e vedere diminuire il proprio reddito.

Un altro concetto importante è la \textbf{mobilità intergenerazionale}, che confronta la posizione economica dei figli adulti con quella dei loro genitori. Se il reddito raggiunto dai figli dipende fortemente dal reddito della famiglia di origine, la mobilità intergenerazionale è bassa. In una situazione simile, nascere in una famiglia ricca o povera influenza molto le possibilità future.

Se invece anche i figli delle famiglie con redditi bassi possono raggiungere posizioni economiche differenti, la mobilità è maggiore. Ciò può indicare che l'istruzione, il mercato del lavoro e le altre istituzioni permettono a un numero più ampio di persone di utilizzare le proprie capacità.

La mobilità intergenerazionale è quindi un indicatore utile dell'\textbf{uguaglianza delle opportunità}. Non misura se tutti ottengono lo stesso risultato, ma aiuta a capire quanto le condizioni familiari determinino il futuro economico di una persona.

Per valutare seriamente la mobilità non basta osservare un singolo caso. La storia di una persona che ha raggiunto il successo partendo da una famiglia povera dimostra che un cambiamento è possibile, ma non ci dice quanto sia probabile per il resto della popolazione. È necessario studiare molte famiglie e seguire i loro redditi per diversi anni.

La mobilità può essere favorita da scuole accessibili e di buona qualità, servizi sanitari, trasporti, possibilità di ottenere credito e un mercato del lavoro aperto. Può invece essere ostacolata dalla discriminazione, dalla forte differenza nella qualità dell'istruzione o dall'assenza di opportunità lavorative in alcune aree.

Una società può quindi essere analizzata attraverso due domande differenti: quanto sono distanti tra loro le condizioni economiche delle persone? E quanto è possibile cambiare la posizione ricevuta alla nascita? Queste domande introducono un problema centrale dell'economia: il rapporto tra la produzione della ricchezza e la sua distribuzione. Nel prossimo capitolo vedremo come capitalismo, socialismo e comunismo propongano risposte differenti a questo problema.

\attivita
\begin{enumerate}
\item Qual è la differenza tra media e mediana?
\item Perché due gruppi con lo stesso reddito medio possono avere distribuzioni molto diverse?
\item Qual è la differenza tra reddito da lavoro e reddito da capitale?
\item Quali fattori possono creare differenze tra i redditi da lavoro?
\item Perché il possesso di un patrimonio può far aumentare le disuguaglianze?
\item Che cosa misura la mobilità intergenerazionale?
\item Perché un singolo caso di successo non dimostra che una società offre le stesse opportunità a tutti?
\end{enumerate}

\lessico
\voce{distribuzione}{modo in cui una quantità è divisa tra più persone}{توزيع}
\voce{media}{somma dei valori divisa per il loro numero}{المتوسط الحسابي}
\voce{mediana}{valore centrale di una serie ordinata}{الوسيط}
\voce{reddito da lavoro}{compenso ottenuto attraverso il lavoro}{الدخل من العمل}
\voce{reddito da capitale}{reddito ottenuto dal possesso di una risorsa}{الدخل من رأس المال}
\voce{patrimonio}{insieme delle risorse economiche possedute}{الثروة}
\voce{disuguaglianza}{differenza nelle condizioni economiche}{عدم المساواة}
\voce{opportunità}{possibilità di raggiungere un risultato}{فرصة}
\voce{mobilità intergenerazionale}{cambiamento di posizione economica tra genitori e figli}{الحراك بين الأجيال}
\voce{eredità}{beni ricevuti dopo la morte di un familiare}{ميراث}

\lingua
Completa con \emph{media, mediana, patrimonio, interesse, mobilità intergenerazionale}: il valore centrale di una serie ordinata è la \ldots; case, terreni e risparmi possono formare il \ldots{} di una persona; il compenso ricevuto per aver prestato denaro è l'\ldots; il confronto tra genitori e figli studia la \ldots.

\scrittura
In circa 150 parole, rispondi alla domanda: una società offre buone opportunità quando tutti ottengono lo stesso reddito oppure quando tutti possono migliorare la propria condizione? Utilizza almeno tre parole del glossario.

\end{lettura}
% FINE LETTURA 08

% =================================================================
% INIZIO LETTURA 09 - circa 1650 parole di testo

\begin{lettura}{1}{Capitalismo, socialismo e sistemi economici}

\subsection*{Produrre ricchezza e distribuirla}
Ogni società deve decidere quali beni produrre, come produrli e come distribuire la ricchezza ottenuta. Un \textbf{sistema economico} è l'insieme delle regole e delle istituzioni attraverso cui una società organizza la produzione, il lavoro, la proprietà e gli scambi. Comprende le decisioni delle famiglie, delle imprese e dello Stato.

Il capitolo precedente ha mostrato che la crescita della ricchezza e la sua distribuzione sono due questioni diverse. Un'economia può produrre di più senza distribuire l'aumento nello stesso modo tra tutti. Da questo problema nascono prospettive differenti. Alcune danno priorità alla libertà economica e all'iniziativa privata; altre attribuiscono maggiore importanza all'uguaglianza e alla proprietà collettiva.

Queste prospettive sono utili per capire le differenze, ma i paesi reali non appartengono quasi mai a un modello puro. Ogni Stato combina mercato, proprietà privata, servizi pubblici e regole in proporzioni diverse. Le combinazioni cambiano anche nel tempo.

\subsection*{Capitalismo e libertà di iniziativa}
Con \textbf{capitalismo} indichiamo un sistema nel quale la proprietà privata dei \textbf{mezzi di produzione}, il lavoro salariato e la produzione per il mercato hanno un ruolo centrale. I mezzi di produzione sono le risorse usate per produrre: per esempio macchinari, edifici, terreni e strumenti. Nel capitalismo questi beni possono appartenere a persone o imprese private.

La \textbf{libertà di iniziativa privata} permette a una persona di aprire un'impresa, scegliere in che cosa investire e provare a introdurre un nuovo prodotto. Chi avvia l'attività assume un rischio: se i ricavi superano i costi ottiene un \textbf{profitto}, mentre se l'impresa non funziona subisce una perdita. Secondo molti economisti, questa possibilità incoraggia l'investimento, l'innovazione e la ricerca di modi più efficienti per produrre.

Il mercato coordina molte di queste decisioni. I prezzi e le vendite segnalano quali prodotti sono richiesti e quali risorse sono più difficili da trovare. Tuttavia, il mercato risponde soprattutto alla domanda di chi desidera e può pagare. Per questa ragione un bisogno sociale importante non si trasforma sempre in una domanda sufficiente a finanziare un servizio.

Il \textbf{liberismo} è la prospettiva che difende un'ampia libertà economica e un intervento limitato dello Stato. Capitalismo e liberismo sono quindi collegati, ma non sono sinonimi: il capitalismo descrive alcune caratteristiche del sistema; il liberismo esprime una posizione su quanto spazio lasciare al mercato e all'iniziativa privata.

\subsection*{Marx e la critica delle disuguaglianze}
\textbf{Karl Marx} fu un filosofo ed economista tedesco dell'Ottocento. Studiò il capitalismo industriale e, insieme a Friedrich Engels, elaborò una delle sue critiche più influenti. Per Marx la divisione decisiva era tra chi possedeva i mezzi di produzione e chi, non possedendoli, doveva vendere il proprio lavoro in cambio di un salario. Chiamò queste due classi \textbf{borghesia} e \textbf{proletariato}.

Nella sua analisi, il lavoratore produce un valore maggiore del salario che riceve. La parte rimanente contribuisce al profitto del proprietario. Se una parte dei profitti viene reinvestita, il proprietario può acquistare altri mezzi di produzione e accrescere ancora il proprio capitale. Per Marx questo meccanismo tende a concentrare ricchezza e potere economico, facendo crescere le disuguaglianze tra proprietari e lavoratori e alimentando il conflitto tra le classi.\fonte{marx}

Questa è una teoria e non una descrizione automatica di ogni paese in ogni periodo. Salari, imposte, sindacati, servizi pubblici e regole possono modificare la distribuzione della ricchezza. La domanda sollevata da Marx rimane però centrale: se il capitale permette di ottenere altro reddito, che cosa impedisce alle differenze iniziali di diventare sempre più grandi?

\subsection*{Socialismo e comunismo}
Il \textbf{socialismo} nasce come risposta alle disuguaglianze e ai rapporti di potere prodotti dal capitalismo industriale. Il termine comprende tradizioni diverse, ma in generale indica l'idea che la produzione debba essere orientata maggiormente verso l'interesse collettivo e che lo Stato, le comunità o i lavoratori debbano controllare una parte importante delle risorse economiche. L'obiettivo è ridurre le disuguaglianze e garantire un accesso più ampio a beni come istruzione, salute e abitazione.

In un sistema socialista possono esistere proprietà pubblica, cooperative e imprese private. Alcune correnti socialiste vogliono correggere il capitalismo attraverso imposte, servizi pubblici e diritti dei lavoratori. Altre propongono di sostituire la proprietà privata dei principali mezzi di produzione con una proprietà collettiva.

In questo secondo percorso, il \textbf{comunismo} rappresenta la forma più estrema: una società senza classi, nella quale i principali mezzi di produzione non appartengono a proprietari privati. Nella teoria di Marx il comunismo avrebbe dovuto superare sia la divisione tra capitalisti e lavoratori sia, nel lungo periodo, la necessità dello Stato come strumento di dominio di una classe sull'altra. La proprietà collettiva dei mezzi di produzione non significa necessariamente eliminare ogni oggetto personale: la distinzione riguarda soprattutto fabbriche, terreni e grandi risorse produttive.

È importante distinguere la teoria dalle esperienze storiche. Gli Stati che nel Novecento si sono definiti comunisti hanno mantenuto governi e apparati pubblici molto forti. Perciò la parola ``comunismo'' può indicare sia un ideale teorico sia sistemi politici realmente esistiti, che non hanno coinciso perfettamente con quell'ideale.

\subsection*{Un continuum tra modelli}
Possiamo immaginare un \textbf{continuum}. A un'estremità si trova un modello nel quale quasi tutte le decisioni economiche e le risorse produttive sono private; all'altra un modello nel quale le principali risorse sono collettive e la produzione è organizzata attraverso decisioni pubbliche. Tra le due estremità esistono molte combinazioni.

Le economie contemporanee sono generalmente \textbf{economie miste}. Possono proteggere la proprietà privata e la libertà d'impresa, ma anche finanziare scuole, ospedali e pensioni, stabilire salari minimi o possedere imprese. La posizione di un paese sul continuum dipende da leggi, storia, rapporti politici e cultura. Inoltre, lo stesso paese può essere più vicino al mercato in un settore e più vicino alla gestione pubblica in un altro.

Anche \textbf{John Maynard Keynes}, economista britannico vissuto tra il 1883 e il 1946, contribuì a cambiare il modo di pensare il rapporto tra Stato e mercato. Dopo aver osservato le crisi economiche e la disoccupazione della prima metà del Novecento, Keynes sostenne che un'economia capitalistica non torna sempre rapidamente alla piena occupazione. Quando famiglie e imprese riducono molto la spesa, lo Stato può sostenere la domanda e limitare la crisi.\fonte{keynes}

Keynes non proponeva di eliminare il capitalismo o la proprietà privata. La sua idea centrale era che lo Stato potesse stabilizzare un'economia di mercato nei momenti di difficoltà. Per questo una politica di intervento pubblico non è automaticamente socialista, così come un'economia capitalistica non è automaticamente liberista.

\subsection*{Stati Uniti ed Europa: il peso della cultura}
Le differenze tra sistemi economici dipendono anche dai valori che una società considera importanti. Negli Stati Uniti la cultura politica ha spesso dato grande spazio alla responsabilità individuale, all'iniziativa privata e alla libertà di scelta. In molti paesi europei hanno avuto maggiore sviluppo l'idea di protezione sociale e la convinzione che alcuni rischi, come malattia, disoccupazione e vecchiaia, debbano essere condivisi attraverso lo Stato.

Questo confronto non deve diventare uno stereotipo. Gli Stati Uniti hanno programmi pubblici e regole economiche; i paesi europei hanno imprese private, mercati e differenze importanti tra loro. Entrambi sono prevalentemente sistemi capitalistici misti. Cambiano però il livello delle imposte, l'estensione dei servizi pubblici e il modo in cui vengono divise le responsabilità tra individuo, impresa e Stato.

\subsection*{Il caso della Cina}
La Cina mostra perché una semplice etichetta può essere insufficiente. La Repubblica Popolare Cinese è guidata dal Partito comunista e definisce il proprio sistema una \textbf{economia socialista di mercato}. Lo Stato possiede o controlla imprese importanti, soprattutto in settori considerati strategici. Allo stesso tempo, le riforme avviate dalla fine degli anni Settanta hanno ampliato lo spazio delle imprese private, degli investimenti e degli scambi di mercato.

La Costituzione cinese rende visibile questa combinazione. L'articolo 6 afferma che la proprietà pubblica è la base del sistema, ma ammette lo sviluppo comune di diverse forme di proprietà. L'articolo 11 definisce le attività individuali e private una componente importante dell'economia socialista di mercato. L'articolo 13 protegge la proprietà privata legittima dei cittadini. Infine, l'articolo 15 stabilisce che lo Stato pratica un'economia socialista di mercato.\fonte{costituzionecinese}

Queste norme non collocano semplicemente la Cina a metà tra capitalismo e comunismo. Mostrano piuttosto la presenza contemporanea di logiche differenti: il mercato coordina molte attività e permette il profitto privato, mentre lo Stato conserva un ruolo molto forte nella proprietà, nel credito, nella pianificazione e negli obiettivi di lungo periodo. Per descrivere il sistema cinese dobbiamo quindi chiederci non solo se esistono imprese private, ma anche chi controlla i settori principali, quali limiti incontra la proprietà e quale autorità stabilisce le priorità economiche.

Non esiste dunque una formula unica per organizzare un'economia. Capitalismo, socialismo e comunismo indicano risposte diverse alle domande sulla proprietà, sulla produzione e sulla distribuzione. Osservare un paese significa capire quali elementi combina e quali obiettivi privilegia: aumentare la ricchezza complessiva, limitarne la concentrazione, proteggere la libertà economica oppure garantire alcuni beni a tutti.

\attivita
\begin{enumerate}
\item Qual è la differenza tra capitalismo e liberismo?
\item Perché la libertà di iniziativa può favorire lo sviluppo economico?
\item Quale meccanismo, secondo Marx, tende ad accrescere le disuguaglianze?
\item Qual è la differenza tra socialismo e comunismo?
\item Perché i sistemi economici reali possono essere rappresentati come un continuum?
\item Chi era Keynes e quale problema cercava di affrontare?
\item Quale differenza generale viene presentata tra Stati Uniti ed Europa?
\item Quali elementi socialisti e quali elementi di mercato convivono in Cina?
\end{enumerate}

\lessico
\voce{sistema economico}{insieme di regole e istituzioni che organizzano produzione e distribuzione}{نظام اقتصادي}
\voce{capitalismo}{sistema fondato principalmente sulla proprietà privata e sulla produzione per il mercato}{الرأسمالية}
\voce{liberismo}{orientamento favorevole alla libertà economica e a un intervento pubblico limitato}{الليبرالية الاقتصادية}
\voce{mezzi di produzione}{risorse e strumenti utilizzati per produrre}{وسائل الإنتاج}
\voce{profitto}{differenza positiva tra ricavi e costi}{ربح}
\voce{socialismo}{insieme di idee che attribuiscono un ruolo maggiore alla proprietà collettiva e all'uguaglianza}{الاشتراكية}
\voce{comunismo}{ideale di una società senza classi e senza proprietà privata dei principali mezzi di produzione}{الشيوعية}
\voce{proprietà collettiva}{proprietà attribuita allo Stato, a una comunità o ai lavoratori}{ملكية جماعية}
\voce{economia mista}{sistema nel quale convivono mercato e intervento pubblico}{اقتصاد مختلط}
\voce{continuum}{insieme di posizioni intermedie tra due estremi}{تدرّج}

\lingua
Completa con \emph{capitalismo, liberismo, socialismo, comunismo, economia mista}: il \ldots{} descrive un sistema basato soprattutto sulla proprietà privata; il \ldots{} è una prospettiva favorevole a un ruolo limitato dello Stato; in una \ldots{} convivono iniziativa privata e intervento pubblico.

\scrittura
In 220--260 parole colloca un paese immaginario lungo il continuum tra capitalismo e socialismo. Descrivi chi possiede le imprese, quali servizi offre lo Stato e come viene affrontato il problema delle disuguaglianze.

\end{lettura}
% FINE LETTURA 09

% =================================================================


\clearpage\phantomsection\addcontentsline{toc}{part}{Terzo trimestre --- Comprendere fenomeni economici più complessi}


% =================================================================
% INIZIO LETTURA 10 - circa 1100 parole di testo

\begin{lettura}{1}{Inflazione e potere d’acquisto}

\subsection*{Con la stessa somma compriamo meno}
Immaginiamo che una confezione di riso costi 20 lire egiziane. Con 100 lire possiamo acquistare cinque confezioni:
\[
100 : 20 = 5.
\]
Dopo un anno, la stessa confezione costa 25 lire. Con le nostre 100 lire possiamo acquistarne soltanto quattro:
\[
100 : 25 = 4.
\]

La somma di denaro non è cambiata, ma permette di comprare una quantità minore di beni. Diciamo quindi che è diminuito il suo \textbf{potere d'acquisto}.

L'\textbf{inflazione} è un aumento generale dei prezzi nel corso del tempo. Nell'esempio è aumentato il prezzo di un solo prodotto. Per parlare di inflazione, invece, dobbiamo osservare molti beni e servizi: alimenti, trasporti, abitazioni, energia, vestiti e altri consumi.

Gli istituti di statistica raccolgono questi prodotti in un \textbf{paniere}. Il costo del paniere viene confrontato in momenti diversi. Se il suo costo complessivo aumenta, il livello generale dei prezzi è aumentato.\fonte{inflazione}

Non tutti i prezzi devono crescere nella stessa misura. Il prezzo del riso può aumentare molto, quello di un vestito può rimanere stabile e quello di un telefono può diminuire. L'inflazione descrive il movimento medio di un insieme di prezzi, non il cambiamento di un singolo prodotto.

\subsection*{Perché aumentano i prezzi?}
Non esiste una sola causa dell'inflazione. In modo semplificato possiamo distinguere tre meccanismi principali:
\begin{enumerate}
\item aumenta rapidamente la domanda;
\item diminuisce l'offerta oppure aumentano i costi di produzione;
\item la quantità di moneta cresce molto più rapidamente della produzione.
\end{enumerate}

Questi meccanismi possono presentarsi contemporaneamente. Un episodio storico può quindi avere più cause.

\subsection*{Prima causa: uno shock della domanda}
La \textbf{domanda} indica la quantità di beni e servizi che famiglie, imprese e Stato vogliono acquistare. Uno \textbf{shock della domanda} è un suo aumento rapido e inatteso.

Se la domanda cresce, le imprese possono provare ad aumentare la produzione. Questo richiede però lavoratori, materie prime, macchine e tempo. Se la produzione non riesce a crescere con la stessa velocità, molte persone cercano di acquistare una quantità limitata di prodotti. I prezzi possono quindi aumentare.

Un esempio è la ripresa successiva alla pandemia di Covid-19. Durante le restrizioni molte persone avevano ridotto alcuni consumi e accumulato risparmi. Alla riapertura delle attività, la domanda di numerosi beni e servizi aumentò rapidamente.

La produzione e i trasporti non riuscirono sempre ad adattarsi con la stessa velocità. Per questa ragione l'inflazione successiva alla pandemia non fu soltanto uno shock della domanda: alla ripresa degli acquisti si aggiunsero problemi nelle catene di produzione e di trasporto.\fonte{postcovid}

Una maggiore domanda non produce sempre inflazione. Se un paese dispone di fabbriche ferme e di molti lavoratori disoccupati, le imprese possono rispondere soprattutto aumentando la produzione. Il rischio di aumento dei prezzi è maggiore quando l'economia non riesce a produrre rapidamente i beni richiesti.

\subsection*{Seconda causa: uno shock dell'offerta}
L'\textbf{offerta} indica la quantità di beni e servizi che le imprese possono produrre e vendere. Uno \textbf{shock dell'offerta} si verifica quando diventa improvvisamente più difficile o più costoso produrre.

Pensiamo all'energia. Il petrolio e il gas vengono utilizzati per produrre elettricità, alimentare le fabbriche e trasportare le merci. Se il loro prezzo aumenta, crescono i costi di molte imprese. Una parte di questi costi può essere trasferita sui prezzi pagati dai consumatori.

Il conflitto che ha coinvolto l'Iran nel 2026 offre un esempio recente. L'Iran si trova in una regione importante per la produzione e il trasporto di petrolio e gas. Il conflitto e il rischio di interruzioni nei trasporti hanno contribuito all'aumento e all'instabilità dei prezzi dell'energia.\fonte{iran2026}

Un'impresa europea o egiziana non deve trovarsi nella zona del conflitto per subirne gli effetti. Se paga di più il carburante, l'elettricità o il trasporto delle merci, i suoi costi aumentano. Lo shock si può così trasmettere da un mercato a molti altri settori.

Altri shock dell'offerta possono essere causati da un raccolto insufficiente, da una siccità, dalla chiusura di una via commerciale o dalla mancanza di una materia prima. In tutti questi casi la quantità disponibile diminuisce oppure il costo di produzione aumenta.

\subsection*{Terza causa: la disponibilità di moneta}
La moneta permette alle persone di acquistare beni e servizi. Se la quantità di moneta e di credito cresce, famiglie e imprese possono aumentare la propria spesa.

Questo non provoca automaticamente inflazione. Se anche la produzione aumenta, la maggiore quantità di moneta può accompagnare una crescita degli scambi. Il problema nasce quando la moneta cresce molto più rapidamente dei beni e dei servizi disponibili.

La banca centrale influenza la quantità di moneta e le condizioni del credito. In alcune situazioni uno Stato può finanziare una parte molto grande delle proprie spese attraverso la creazione di nuova moneta. Se questo processo continua mentre la produzione non cresce, i prezzi possono aumentare rapidamente.

Un caso storico estremo si verificò in Bolivia nella prima metà degli anni Ottanta. Lo Stato aveva grandi disavanzi: le sue spese erano molto superiori alle entrate. Una parte di questi disavanzi venne finanziata attraverso una forte espansione monetaria. Nel 1985 l'aumento dei prezzi nell'arco dell'anno superò l'11.000 per cento.\fonte{bolivia}

Quando i prezzi aumentano in modo estremamente rapido e la moneta perde velocemente valore si parla di \textbf{iperinflazione}. Le persone cercano di spendere il denaro appena lo ricevono, perché temono che dopo poco tempo permetterà di comprare molto meno. Questo comportamento può accelerare ulteriormente l'aumento dei prezzi.

L'iperinflazione è un caso raro ed estremo. Non dobbiamo confonderla con ogni normale aumento dell'inflazione. L'esempio boliviano mostra però che una crescita molto forte e prolungata della moneta, senza un corrispondente aumento della produzione, può distruggere la fiducia nella valuta.

\begin{center}
\begin{minipage}{\linewidth}\centering\small
\begin{tabularx}{\linewidth}{@{}lXX@{}}
\toprule
\textbf{Causa} & \textbf{Meccanismo} & \textbf{Esempio} \\
\midrule
Domanda & Gli acquisti crescono più velocemente della produzione. & Ripresa dopo il Covid-19. \\
Offerta & Diminuisce la produzione oppure aumentano i costi. & Energia e conflitto iraniano. \\
Moneta & La moneta cresce molto più rapidamente dei beni disponibili. & Bolivia nel 1985. \\
\bottomrule
\end{tabularx}
\captionof{table}{Tre cause principali dell'inflazione.}
\end{minipage}
\end{center}

\subsection*{Tasso di cambio e prodotti importati}
Un paese può consumare molti prodotti realizzati all'estero. Per acquistarli deve cambiare la propria moneta con la valuta richiesta dal venditore.

Il \textbf{tasso di cambio} indica il prezzo di una moneta espresso in un'altra moneta. Immaginiamo che un componente importato costi 10 dollari.

Se un dollaro vale 30 lire egiziane, il componente costa:
\[
10 \times 30 = 300 \text{ LE}.
\]
Se successivamente un dollaro vale 40 lire egiziane, lo stesso componente costa:
\[
10 \times 40 = 400 \text{ LE}.
\]

Il produttore straniero non ha modificato il prezzo in dollari. Il prodotto è diventato più costoso in Egitto perché sono necessarie più lire per acquistare un dollaro.

Quando una moneta perde valore rispetto alle valute straniere parliamo di \textbf{deprezzamento} o, in alcuni sistemi monetari, di svalutazione. Se un paese importa una parte importante degli alimenti, dell'energia o dei componenti industriali che utilizza, il deprezzamento può contribuire all'inflazione.

L'effetto è maggiore quando è difficile sostituire i prodotti importati con prodotti nazionali. Se le imprese riescono a utilizzare materie prime locali, il legame tra cambio e prezzi può essere più debole.

\subsection*{Inflazione, stipendi e potere d'acquisto}
Per comprendere gli effetti dell'inflazione dobbiamo confrontare l'aumento dei prezzi con l'aumento dei redditi.

Supponiamo che lo stipendio di una lavoratrice aumenti del 5 per cento, mentre i prezzi aumentano mediamente del 10 per cento. La lavoratrice riceve più denaro, ma il suo stipendio non cresce abbastanza per acquistare gli stessi beni di prima. Il suo reddito \textbf{nominale} è aumentato, mentre il suo reddito \textbf{reale}, cioè il suo potere d'acquisto, è diminuito.

L'inflazione colpisce in modo particolare chi riceve un reddito che cambia lentamente. Per sapere se una persona è diventata economicamente più ricca non basta quindi osservare quanti soldi riceve: dobbiamo anche osservare che cosa può comprare con quei soldi.

Le cause dell'inflazione possono essere diverse, ma la domanda finale rimane la stessa: i redditi riescono a crescere quanto il costo della vita?

\attivita
\begin{enumerate}
\item Quante confezioni di riso si possono acquistare con 100 LE prima e dopo l'aumento del prezzo?
\item Perché l'aumento del prezzo di un solo prodotto non basta per parlare di inflazione?
\item Che cosa accade quando la domanda cresce più rapidamente della produzione?
\item Perché un aumento del prezzo dell'energia può influenzare molti altri prodotti?
\item In quale situazione una forte crescita della moneta può provocare inflazione?
\item Che cosa accadde in Bolivia nel 1985?
\item Spiega il passaggio dal tasso di cambio al prezzo del componente importato.
\item Quando un aumento dello stipendio non produce un aumento del potere d'acquisto?
\end{enumerate}

\lessico
\voce{inflazione}{aumento generale dei prezzi nel tempo}{تضخم}
\voce{potere d'acquisto}{quantità di beni e servizi acquistabile con una somma}{قوة شرائية}
\voce{paniere}{insieme di beni e servizi utilizzato per osservare i prezzi}{سلة}
\voce{domanda}{quantità che i consumatori vogliono acquistare}{طلب}
\voce{offerta}{quantità che le imprese possono produrre e vendere}{عرض}
\voce{shock}{cambiamento improvviso delle condizioni economiche}{صدمة}
\voce{iperinflazione}{aumento estremamente rapido dei prezzi}{تضخم مفرط}
\voce{tasso di cambio}{prezzo di una moneta espresso in un'altra}{سعر الصرف}
\voce{deprezzamento}{perdita di valore di una moneta rispetto a un'altra}{انخفاض قيمة العملة}
\voce{reddito reale}{reddito considerato rispetto ai prezzi}{الدخل الحقيقي}

\lingua
Completa le frasi con \emph{domanda, offerta, potere d'acquisto, tasso di cambio}: se i prezzi aumentano e lo stipendio resta uguale, diminuisce il \ldots; una siccità può ridurre l'\ldots{} di prodotti agricoli; alla riapertura delle attività aumentò la \ldots; il prezzo di una valuta espresso in un'altra è il \ldots.

\scrittura
Scrivi 180--200 parole per spiegare due possibili cause dell'inflazione. Utilizza un esempio storico e concludi spiegando l'effetto sul potere d'acquisto di una famiglia.

\end{lettura}
% FINE LETTURA 10

% =================================================================
% INIZIO LETTURA 11 - circa 1821 parole di testo

\begin{lettura}{1}{Le banche centrali e i tassi di interesse}


\subsection*{Un'istituzione diversa da una banca commerciale}
Quando una famiglia apre un conto o un'impresa chiede un prestito, si rivolge generalmente a una banca commerciale. Una \textbf{banca centrale} svolge funzioni differenti: opera al centro del sistema monetario e interviene sulle condizioni in cui circolano moneta e credito. Non è semplicemente una banca commerciale più grande e non decide direttamente ogni prezzo dell'economia.

Per comprendere il suo ruolo dobbiamo collegare tre letture precedenti. Abbiamo visto che le banche offrono pagamenti e credito, che i tassi influenzano il costo dei finanziamenti e che l'inflazione riguarda il livello generale dei prezzi. La banca centrale agisce su questi collegamenti attraverso strumenti specifici. Gli effetti passano per le decisioni di banche, famiglie e imprese.

Le funzioni e gli obiettivi sono definiti dal quadro istituzionale. Non tutte le banche centrali hanno esattamente lo stesso mandato. Per sapere che cosa deve fare una particolare istituzione bisogna leggere i suoi documenti. Il nome «banca centrale» identifica una funzione generale, ma non rende identiche tutte le regole nazionali.

\subsection*{Prezzi stabili non significa prezzi immobili}
La \textbf{stabilità dei prezzi} riguarda un andamento contenuto e prevedibile del livello generale dei prezzi. Non significa fissare il prezzo di ogni prodotto. In un'economia i prezzi relativi possono cambiare: un bene può diventare più costoso perché è più richiesto, mentre un altro può costare meno grazie a un'innovazione.

Nell'area euro, la Banca centrale europea mira a un'inflazione del due per cento nel medio termine. L'obiettivo riguarda l'indice generale e non ogni singolo prezzo. Il riferimento al medio termine riconosce che gli shock e gli effetti delle decisioni richiedono tempo. Il due per cento non è un tasso di interesse e non è una previsione valida automaticamente per ogni anno.\fonte{bceobiettivo}

Anche la Banca centrale d'Egitto indica tra i propri obiettivi la stabilità dei prezzi e la solidità del sistema monetario e bancario, nel quadro della politica economica generale dello Stato. Per conoscere obiettivi numerici e decisioni di un particolare periodo occorre consultare i documenti riferiti a quel periodo.\fonte{cbe}

\subsection*{Moneta della banca centrale e depositi bancari}
Abbiamo imparato che esistono banconote e depositi. Le banche commerciali utilizzano inoltre conti presso la banca centrale per regolare pagamenti tra loro. I saldi di questi conti sono chiamati \textbf{riserve}. Non sono la stessa cosa del capitale della banca e non coincidono con tutti i risparmi delle famiglie.

Immaginiamo che una cliente della banca A paghi un fornitore con un conto presso la banca B. Il pagamento modifica i depositi dei clienti e richiede un regolamento tra le due banche. Questa dimensione aiuta a comprendere perché una banca deve gestire la propria liquidità anche quando concede prestiti tramite accrediti contabili.

Le banche non possono ignorare perdite, pagamenti e regole. Un prestito che non viene restituito può ridurre le risorse capaci di assorbire perdite; un pagamento verso un'altra banca richiede liquidità. Sono problemi collegati, ma differenti. La banca centrale opera in questo sistema attraverso le condizioni applicate alle proprie operazioni e altri strumenti previsti dal suo mandato.

\subsection*{I tassi di politica monetaria}
Un \textbf{tasso di politica monetaria} è un tasso fissato dalla banca centrale per determinate operazioni. Può riguardare, per esempio, depositi o finanziamenti delle banche presso di essa. Questi tassi influenzano altri tassi a breve termine e, attraverso vari passaggi, le condizioni del credito nell'economia.\fonte{bcemonetaria}

Il tasso applicato a un prestito a un'impresa non coincide automaticamente con quello della banca centrale. Include valutazioni sulla durata, sul rischio del debitore, sui costi e sulle condizioni di concorrenza. Due imprese possono quindi ricevere offerte diverse anche nello stesso giorno e nello stesso paese.

Per questo, quando una notizia dice che «i tassi aumentano», dobbiamo chiedere quali tassi. Può riferirsi a una decisione ufficiale, ai prestiti per le famiglie oppure al rendimento dei titoli di Stato. Le serie possono muoversi in modo collegato, ma hanno significati diversi e possono reagire con tempi differenti.

\subsection*{Dal tasso al prestito}
Supponiamo che il costo di un nuovo finanziamento aumenti. Un'impresa che valutava l'acquisto di una macchina deve ricalcolare il progetto. Se i ricavi attesi restano invariati e il costo del debito cresce, alcuni investimenti possono diventare meno convenienti. L'impresa può rinviare l'acquisto, ridurre il progetto o cercare altre risorse.

Anche una famiglia può rinviare una spesa finanziata a debito. Il passaggio non è però uguale per tutti i contratti. Chi ha un prestito a tasso fisso può mantenere le condizioni già stabilite; chi ha un tasso variabile può subire una modifica secondo le clausole previste. I nuovi prestiti e quelli già esistenti vanno quindi distinti.

Un aumento dei tassi può rendere più interessante conservare alcune forme di risparmio remunerato. Tuttavia, la scelta dipende dalle necessità della famiglia, dal rischio e dall'inflazione attesa. Non tutte le persone possono ridurre i consumi: alcune destinano quasi tutto il reddito alle spese essenziali.

\subsection*{Dal credito alla domanda}
Se molte decisioni di spesa vengono ridotte o rinviate, la \textbf{domanda complessiva} può rallentare. Le imprese possono ricevere meno ordini e trovare più difficile aumentare i prezzi. Questa è una parte del meccanismo attraverso cui una politica monetaria più restrittiva può ridurre le pressioni inflazionistiche.

La parola \textbf{restrittiva} indica un orientamento che rende le condizioni monetarie e finanziarie meno favorevoli all'espansione della spesa. Una politica \textbf{espansiva} cerca invece di sostenerle. Il significato concreto dipende dal contesto: lo stesso tasso nominale può risultare più o meno restrittivo a seconda dell'inflazione attesa e delle altre condizioni economiche.

Se la domanda rallenta, possono diminuire anche produzione e opportunità di lavoro. Il controllo dell'inflazione può quindi comportare costi nel breve periodo. Una banca centrale deve considerare intensità, durata e origine delle pressioni sui prezzi. Non dispone di uno strumento capace di aumentare automaticamente ogni beneficio e ridurre ogni costo.

\subsection*{Aspettative e comunicazione}
Le decisioni economiche guardano al futuro. Un'impresa che firma un contratto pluriennale considera costi e prezzi attesi. Una famiglia che assume un debito considera entrate e rate future. Per questo la \textbf{comunicazione} della banca centrale può influenzare le aspettative, oltre all'effetto immediato dei tassi.

Se l'istituzione spiega obiettivi, dati e ragioni delle decisioni, gli operatori possono comprendere meglio come potrebbe reagire a condizioni nuove. Questo non significa promettere una sequenza immutabile di tassi. Una decisione può cambiare quando cambiano le informazioni. La credibilità riguarda anche la coerenza tra mandato, spiegazioni e comportamento.

Una previsione è diversa da un impegno. «Prevediamo una riduzione dell'inflazione» descrive uno scenario atteso; «abbiamo deciso un nuovo tasso» descrive un'azione. Nel leggere un comunicato dobbiamo distinguere queste parti. Le parole «se», «potrebbe» e «nel caso in cui» indicano condizioni, non semplici esitazioni linguistiche.

\subsection*{Il canale del cambio}
I tassi possono influenzare anche i movimenti di risorse finanziarie tra paesi e il cambio. Un rendimento più elevato può attirare investitori, ma contano pure rischio, aspettative e possibilità di trasferire le risorse. Non esiste una relazione certa in cui ogni aumento dei tassi rafforzi sempre la moneta.

Se una moneta si rafforza, alcuni beni importati possono diventare meno costosi in moneta nazionale. Se si indebolisce, può avvenire il contrario. Il passaggio ai prezzi al consumo dipende dalla quota importata, dai contratti e dalle decisioni delle imprese. Ritroviamo così il collegamento tra cambio e inflazione studiato nella lettura precedente.

Un'impresa esportatrice può essere coinvolta in modo diverso da un'impresa che importa molti componenti. La stessa variazione del cambio modifica ricavi e costi in modi differenti. Per valutare l'effetto complessivo bisogna quindi conoscere la struttura delle attività e non soltanto la direzione del movimento della moneta.

\subsection*{Che cosa la banca centrale non produce}
Una banca centrale non può fabbricare direttamente gas, grano o componenti elettronici. Se uno shock riduce la disponibilità di un bene essenziale, modificare i tassi non elimina materialmente quella scarsità. Può però influenzare la domanda, le aspettative e il modo in cui lo shock si trasmette al resto dell'economia.

Per affrontare problemi produttivi possono essere necessari altri strumenti: riparare infrastrutture, diversificare fornitori, migliorare trasporti o sostenere capacità tecniche. Questi interventi appartengono a politiche differenti e hanno tempi diversi. La distinzione aiuta a comprendere perché la politica monetaria non sia l'unica politica economica.

Anche le decisioni fiscali, cioè su entrate e spese pubbliche, possono influenzare domanda e prezzi. Una riduzione di un'imposta e una modifica di un tasso sono operazioni differenti. Possono avere effetti che si rafforzano oppure si compensano. Per leggere il quadro complessivo dobbiamo riconoscere quale istituzione utilizza quale strumento.

\subsection*{I tempi della trasmissione}
La \textbf{trasmissione monetaria} è il percorso attraverso cui una decisione influenza condizioni finanziarie, spesa e prezzi. Non avviene tutta nello stesso giorno. Alcuni mercati reagiscono rapidamente; contratti e progetti già avviati possono modificarsi più lentamente. Il ritardo rende più difficile valutare una decisione soltanto dai dati del mese successivo.

Immaginiamo un'impresa che abbia già ordinato un impianto e ottenuto un finanziamento a tasso fisso. Un nuovo aumento dei tassi può non cambiare subito quel progetto. Può però influenzare il prossimo investimento o le decisioni di altri clienti. Gli effetti si accumulano attraverso molte scelte individuali.

Nel frattempo possono intervenire nuovi eventi: un miglioramento dei trasporti, una variazione dei prezzi energetici o una crisi estera. Perciò osservare inflazione e tassi nello stesso periodo non basta a ricostruire una causa unica. La lettura economica deve considerare sequenze, ritardi e condizioni esterne.

\subsection*{Un caso immaginario da ricostruire}
In una piccola economia, molte imprese lavorano quasi al massimo della capacità e i prezzi aumentano rapidamente. La banca centrale alza un proprio tasso. Le banche rivedono gradualmente le condizioni dei nuovi prestiti. Alcune imprese rinviano investimenti e alcune famiglie rimandano acquisti. La domanda rallenta e le pressioni sui prezzi possono ridursi.

Ora cambiamo la situazione iniziale. Una tempesta interrompe l'arrivo di componenti essenziali, mentre la domanda non cresce. I prezzi aumentano perché produrre è più difficile. Alzare i tassi non riapre il porto. Può incidere su altri meccanismi, ma comporta una valutazione diversa rispetto al primo caso. Le due storie mostrano perché l'origine del problema conta.

Questi esempi non forniscono una ricetta per una decisione reale. Servono a leggere il collegamento tra problema, strumento e risultato atteso. Una spiegazione completa deve indicare attraverso quali passaggi dovrebbe funzionare l'intervento e quali condizioni potrebbero limitarlo.

\subsection*{Una decisione comune, effetti differenti}
Immaginiamo due imprese che vendano prodotti simili. Una finanzia gli investimenti soprattutto con risorse accumulate; l'altra dipende da nuovi prestiti. Un aumento del costo del credito può incidere in modo diverso sulle loro decisioni. Anche la durata dei contratti e la situazione dei clienti possono modificare il risultato.

Questa differenza aiuta a leggere il linguaggio delle medie. Una politica può frenare la domanda complessiva senza ridurre ogni singola spesa. Una banca può aumentare un tasso senza modificare nello stesso momento tutte le condizioni. Quando riassumiamo un testo, conserviamo quindi parole come «in generale», «gradualmente» e «può». Non sono decorazioni: delimitano ciò che l'autore sta affermando e impediscono di trasformare una relazione complessiva in una regola individuale assoluta.

\subsection*{Leggere un comunicato}
Quando incontriamo un comunicato, cerchiamo prima la decisione effettiva: quale tasso è stato modificato e da quando? Poi distinguiamo le ragioni, i dati osservati e lo scenario futuro. Infine individuiamo i rischi indicati dall'istituzione. In questo modo un testo tecnico diventa una sequenza di informazioni riconoscibili.

Possiamo riassumerlo con una struttura semplice: «L'istituzione ha deciso… perché osserva…; si attende… ma segnala…». Questa struttura conserva la differenza tra azione, motivazione, previsione e incertezza. È una competenza di lettura utile anche per documenti che riguardano ambiente, investimenti pubblici e cooperazione internazionale.


\attivita
\begin{enumerate}

\item Distingui banca centrale e banca commerciale.

\item Il 2 per cento citato per la BCE è un tasso di interesse?

\item Ricostruisci i passaggi dal tasso al prestito e dal prestito alla domanda.

\item Perché prestiti a tasso fisso e variabile possono reagire diversamente?

\item Quale limite incontra la politica monetaria nel caso del porto chiuso?

\item Perché gli effetti non sono immediati?

\end{enumerate}
\lessico

\voce{mandato}{compiti e obiettivi assegnati a un’istituzione}{تفويض}

\voce{riserve}{saldi delle banche presso la banca centrale}{احتياطيات}

\voce{restrittivo}{che frena l’espansione monetaria e della domanda}{انكماشي}

\voce{espansivo}{che sostiene condizioni favorevoli alla spesa}{توسعي}

\voce{trasmissione}{passaggio degli effetti attraverso più relazioni}{انتقال}

\voce{credibilità}{fiducia nella coerenza di un soggetto}{مصداقية}

\voce{previsione}{descrizione di un esito futuro atteso}{تنبؤ}

\lingua
Dividi un paragrafo in quattro categorie: decisione, motivazione, previsione e condizione. Scrivi poi un riassunto usando «ha deciso», «poiché», «si attende», «se».
\scrittura
Confronta in 250--300 parole i due casi immaginari della lettura. Spiega perché la stessa variazione dei prezzi può richiedere analisi differenti.

\end{lettura}
% FINE LETTURA 11

% =================================================================
% INIZIO LETTURA 12 - circa 2051 parole di testo

\begin{lettura}{1}{Crescita, produttività e innovazione}


\subsection*{Produrre di più nel tempo}
Quando parliamo di \textbf{crescita economica} ci riferiamo a un aumento della produzione, generalmente misurato attraverso il PIL reale. La parola «reale» ricorda che vogliamo distinguere le quantità e la qualità della produzione dai semplici aumenti dei prezzi. Se il valore monetario sale soltanto perché i prezzi aumentano, non possiamo concludere che il paese produca più beni e servizi.

Anche la crescita totale va distinta da quella per abitante. Se produzione e popolazione aumentano nella stessa proporzione, la produzione media per persona resta invariata. Per comprendere l'evoluzione delle possibilità materiali di una società è quindi utile osservare entrambe le dimensioni, senza dimenticare la distribuzione.

Il benessere rimane un concetto più ampio. Una produzione maggiore può offrire risorse per servizi e consumi, ma non garantisce da sola salute, tempo libero o qualità ambientale. La domanda economica non è soltanto «quanto produciamo?», ma anche «come produciamo, che cosa otteniamo e come vengono distribuiti benefici e costi?».

\subsection*{Lavoro, capitale e organizzazione}
Per produrre utilizziamo lavoro, strumenti e risorse naturali. Il \textbf{capitale produttivo} comprende beni durevoli impiegati nella produzione, come macchine, edifici e reti. Il capitale finanziario, invece, riguarda risorse e strumenti di finanziamento. Le due espressioni sono collegate, ma non coincidono.

Un laboratorio può aumentare la produzione assumendo più persone o acquistando più macchine. Può anche migliorare l'organizzazione con le stesse risorse. Se i materiali arrivano in ritardo e gli operai aspettano, una parte del tempo disponibile non viene utilizzata produttivamente. Ridurre l'attesa può aumentare il risultato senza allungare l'orario di lavoro.

La produzione dipende inoltre dalla combinazione delle risorse. Una macchina sofisticata può essere poco utile se nessuno sa utilizzarla o se manca elettricità affidabile. Competenze, infrastrutture e manutenzione sono quindi complementari agli investimenti materiali. La semplice presenza di uno strumento non coincide con la sua utilizzazione efficace.

\subsection*{Che cosa misura la produttività}
La \textbf{produttività} mette in relazione il prodotto ottenuto e le risorse impiegate. Una misura semplice è la produzione per ora di lavoro. Se un laboratorio realizza quaranta oggetti in venti ore, produce in media due oggetti per ora. Se con le stesse venti ore ne realizza sessanta, la produttività del lavoro sale a tre oggetti per ora.

Questo aumento non significa necessariamente che i lavoratori si muovano più velocemente o si impegnino di più. Può dipendere da strumenti migliori, materiali più adatti o procedure più chiare. Attribuire tutta la produttività alla persona ignora il contesto in cui lavora.

Anche il confronto richiede attenzione alla qualità. Produrre più oggetti difettosi non equivale a produrre più beni utilizzabili. Se il risultato comprende prodotti diversi, contare i pezzi può non essere sufficiente. Per confronti più ampi si utilizzano misure di valore reale, pur con difficoltà di misurazione.

\subsection*{Un laboratorio immaginario}
Consideriamo un laboratorio che produce piccoli supporti metallici. In una giornata di otto ore, quattro lavoratori realizzano centoventotto pezzi utilizzabili. Le ore complessive di lavoro sono trentadue: quattro persone per otto ore. La produttività media è quattro pezzi per ora di lavoro.

Il responsabile modifica la disposizione degli strumenti e introduce un dispositivo che riduce gli errori. Con le stesse persone e lo stesso orario si ottengono centosessanta pezzi utilizzabili. La produttività sale a cinque pezzi per ora, con un aumento del venticinque per cento. I dati sono inventati per comprendere la relazione.

Per valutare il progetto dobbiamo però considerare anche il costo del dispositivo, la formazione e la manutenzione. Un aumento della produttività tecnica può essere interessante, ma la convenienza economica dipende da costi e benefici nel tempo. Inoltre, vendere più pezzi richiede clienti: capacità produttiva e domanda non sono sinonimi.

\subsection*{Innovazione di prodotto}
Un'\textbf{innovazione di prodotto} introduce un bene o servizio nuovo oppure migliora in modo significativo uno esistente. Un dispositivo che consuma meno energia e svolge nuove funzioni può rappresentare un'innovazione di prodotto. Anche un servizio di prenotazione più accessibile può modificare il modo in cui viene soddisfatto un bisogno.

L'innovazione non coincide sempre con un oggetto completamente nuovo per il mondo. Un miglioramento può essere importante per gli utenti che lo adottano. Tuttavia, cambiare soltanto il colore di una confezione non ha necessariamente lo stesso significato di introdurre una funzione utile. Per descrivere l'innovazione dobbiamo spiegare che cosa cambia.

Un nuovo prodotto può creare un mercato, sostituire prodotti precedenti oppure affiancarli. La domanda dipende da prezzo, utilità, abitudini e possibilità di accesso. Una tecnologia disponibile non viene automaticamente adottata da tutti. Le persone devono conoscerla, comprenderla e poter sostenere i costi.

\subsection*{Innovazione di processo}
Un'\textbf{innovazione di processo} modifica il modo in cui si produce o si organizza un'attività. Può ridurre tempi, errori e consumo di materiali. Un sistema che individua un difetto prima della fase finale permette, per esempio, di evitare lavorazioni inutili. Il prodotto può restare simile, mentre cambia il percorso necessario a ottenerlo.

Le due forme di innovazione possono combinarsi. Un nuovo programma informatico è un prodotto per chi lo vende e può diventare uno strumento di processo per chi lo utilizza. La classificazione dipende quindi dal punto di vista e dalla funzione osservata.

Anche procedure semplici possono essere innovative per un'organizzazione. Se un'officina registra gli interventi e organizza meglio le scorte, può ridurre gli acquisti urgenti e i tempi di attesa. Non ogni miglioramento richiede una macchina costosa. La conoscenza dell'attività e la collaborazione dei lavoratori possono essere decisive.

\subsection*{Invenzione, adozione e diffusione}
Un'\textbf{invenzione} introduce una soluzione o un'idea tecnica. Per produrre effetti economici deve spesso essere sviluppata, adottata e diffusa. Tra la prima dimostrazione e l'utilizzo quotidiano possono passare anni. Servono prove, adattamenti, finanziamenti e persone capaci di impiegare la soluzione.

La \textbf{diffusione} indica l'estensione dell'utilizzo tra imprese, persone o territori. Un'impresa può imparare da un fornitore, da un lavoratore formato altrove o da una collaborazione. Le conoscenze circolano attraverso relazioni concrete, non soltanto attraverso la disponibilità di un manuale.

La capacità di adottare una tecnologia dipende anche da competenze precedenti. Un programma avanzato può essere difficile da utilizzare se i dati sono incompleti o le procedure non sono organizzate. Per questo la formazione continua e la manutenzione delle capacità possono accompagnare l'acquisto di strumenti.

\subsection*{Automazione e compiti}
L'\textbf{automazione} affida a sistemi tecnici attività prima svolte, in tutto o in parte, dalle persone. Può riguardare movimenti ripetitivi, controlli o elaborazioni di dati. Per comprenderne gli effetti è utile distinguere una professione dall'insieme dei compiti che la compongono.

Un tecnico può leggere misure, individuare problemi, parlare con il cliente e controllare la sicurezza. Un nuovo sistema può automatizzare soltanto una parte di questi compiti. Alcune competenze possono diventare meno richieste, mentre altre acquistano importanza. Non ogni automazione elimina quindi un'intera professione, ma non è neppure corretto escludere possibili perdite di posti di lavoro.

Il risultato dipende da come cambia la domanda, da quali nuove attività emergono e dalla possibilità di formare i lavoratori. Una maggiore efficienza può ridurre i costi e ampliare le vendite; può anche ridurre il lavoro necessario per una produzione invariata. Sono meccanismi diversi che devono essere osservati nel contesto.

\subsection*{Il costo relativo delle risorse}
Un'impresa sceglie tra tecniche produttive considerando i costi e la disponibilità delle risorse. Una macchina può richiedere un investimento iniziale elevato ma ridurre alcune spese successive. Il lavoro umano può offrire flessibilità, capacità di adattamento e competenze non facilmente sostituibili.

Se il costo del finanziamento aumenta, acquistare una macchina può diventare meno conveniente. Se l'energia è inaffidabile, una tecnica dipendente da un'alimentazione continua può presentare rischi. Perciò il rapporto tra lavoro e capitale dipende anche da infrastrutture e condizioni finanziarie, non soltanto dal salario.

Due imprese in paesi differenti possono utilizzare tecniche diverse per realizzare prodotti simili. Questa differenza non dimostra automaticamente che una sia più capace dell'altra. Le condizioni economiche possono rendere conveniente una combinazione diversa. Per confrontarle bisogna considerare costi, qualità, tempi e vincoli.

\subsection*{Istruzione e capacità produttiva}
Le competenze aiutano a utilizzare strumenti, interpretare problemi e adattare procedure. L'istruzione può contribuire alla capacità produttiva, ma il titolo di studio da solo non descrive tutto ciò che una persona sa fare. Contano qualità dell'apprendimento, esperienza e possibilità di applicare le conoscenze.

Anche leggere un manuale tecnico è una competenza economica. Se un lavoratore comprende un'istruzione, può ridurre errori e riconoscere rischi. Il vocabolario specialistico serve a comunicare in modo preciso. Le letture di questo libretto collegano quindi l'apprendimento linguistico alla capacità di comprendere ambienti di lavoro e decisioni collettive.

La formazione può avere benefici che vanno oltre la singola persona. Un lavoratore condivide conoscenze con colleghi e può migliorare l'organizzazione. Tuttavia, chi sostiene i costi della formazione non riceve necessariamente tutti i benefici. Ritroviamo così il tema degli effetti che superano una singola relazione economica.

\subsection*{Infrastrutture e istituzioni}
Un'infrastruttura è una struttura o rete che sostiene molte attività: trasporti, energia, acqua e comunicazioni ne sono esempi. Il suo valore dipende anche dalle connessioni. Un tratto ferroviario isolato può essere meno utile di una rete che collega luoghi di produzione e consumo.

Anche le istituzioni influenzano le possibilità di investimento. Contratti comprensibili, procedure affidabili e tempi prevedibili possono ridurre l'incertezza. Questo non significa che basti una sola regola per produrre crescita. Le condizioni si combinano e possono rafforzarsi o limitarsi reciprocamente.

Un'impresa può avere una buona idea ma non trovare finanziamenti; può ottenere finanziamenti ma non disporre di energia; può produrre bene ma non raggiungere i clienti. Per capire lo sviluppo bisogna seguire l'intera catena di condizioni, evitando di cercare una spiegazione unica per ogni paese.

\subsection*{Ricerca e incertezza}
La \textbf{ricerca} cerca nuove conoscenze e soluzioni. Il risultato non è sempre prevedibile: molti tentativi non producono immediatamente un'applicazione. Un'organizzazione deve quindi sostenere costi prima di sapere se otterrà un beneficio. Questa incertezza distingue la ricerca da un acquisto di risultati già disponibili.

Imprese, università e istituzioni pubbliche possono svolgere ruoli diversi. Una conoscenza può essere condivisa ampiamente oppure protetta attraverso diritti specifici. La protezione può offrire incentivi, mentre la diffusione può ampliare l'utilizzo. Le politiche devono considerare entrambi gli aspetti.

Anche un risultato non previsto può essere utile in un altro contesto. Per questo collaborazione e circolazione delle idee possono avere valore. Non dobbiamo però confondere una possibilità con una garanzia: investire in ricerca richiede scelte, valutazioni e capacità di apprendere anche dagli insuccessi.

\subsection*{Crescita e distribuzione dei benefici}
Se un'impresa produce di più con le stesse risorse, il beneficio può assumere forme diverse: prezzi più bassi, salari maggiori, profitti più alti o nuovi investimenti. La ripartizione dipende da mercato, contratti e decisioni. Non esiste un passaggio automatico da maggiore produttività a benefici identici per tutti.

Anche i territori possono beneficiare in modi diversi dell'innovazione. Un'area con competenze e infrastrutture può attirare attività, mentre un'altra può perdere occupazione. Le politiche di formazione e collegamento possono influenzare queste differenze, ma richiedono tempo e risorse.

Per valutare una trasformazione tecnologica dobbiamo quindi osservare sia il risultato complessivo sia le persone coinvolte. Un aumento della produzione totale può convivere con difficoltà per alcuni gruppi. Riconoscere questa possibilità permette di discutere gli adattamenti necessari senza negare il valore dell'innovazione.

\subsection*{La dimensione ambientale}
Produrre di più può richiedere più energia e materiali, ma le tecnologie possono anche ridurre il consumo per unità di prodotto. Le due variazioni vanno distinte. Se ogni oggetto richiede meno energia ma il numero degli oggetti aumenta molto, il consumo totale può comunque crescere.

Immaginiamo che il consumo per pezzo scenda da dieci a otto unità, mentre la produzione passi da cento a centocinquanta pezzi. Il consumo totale passa da mille a milleduecento unità. L'efficienza migliora, ma il totale aumenta. È un esempio didattico di come intensità e quantità complessiva possano muoversi diversamente.

Questo collegamento prepara la prossima lettura. Lo sviluppo sostenibile richiede di osservare produzione, distribuzione e limiti ambientali insieme. Una tecnologia va descritta attraverso il problema che affronta, le risorse che utilizza e le conseguenze che produce durante il suo impiego.

\subsection*{Il tempo risparmiato}
Un miglioramento produttivo può anche liberare tempo. Se un'attività richiede meno ore, quelle ore possono essere dedicate ad altri compiti, a una produzione maggiore o al riposo. La scelta dipende dall'organizzazione e dagli accordi. Per valutare il beneficio dobbiamo quindi osservare come viene utilizzato il tempo risparmiato, non soltanto contarlo. Lo stesso cambiamento tecnico può tradursi in risultati sociali differenti.

\subsection*{Leggere il lungo periodo}
La crescita di lungo periodo nasce da cambiamenti che si accumulano: strumenti, competenze, organizzazioni e reti. Una decisione può avere effetti dopo molti anni e dipendere da altre decisioni. Perciò una singola variazione annuale non descrive tutto il percorso di sviluppo.

Nel leggere un testo cerchiamo i meccanismi: quale risorsa aumenta? Quale processo migliora? Quale ostacolo viene ridotto? Poi chiediamo come i benefici vengono distribuiti e quali nuovi problemi emergono. Questa sequenza permette di comprendere l'innovazione come un processo economico e sociale, non soltanto come un elenco di oggetti moderni.



\begin{center}\begin{minipage}{\linewidth}\centering\small\begin{tabular}{lrr}\toprule
Laboratorio immaginario & Prima & Dopo\\\midrule
Ore di lavoro totali & 32 & 32\\
Pezzi utilizzabili & 128 & 160\\
Pezzi per ora & 4 & 5\\\bottomrule
\end{tabular}
\captionof{table}{Stesse ore di lavoro, produttività maggiore del 25\%.}
\end{minipage}\end{center}


\attivita
\begin{enumerate}

\item Calcola la produttività del laboratorio prima e dopo la modifica.

\item Perché maggiore produttività non significa necessariamente maggiore sforzo individuale?

\item Distingui innovazione di prodotto e di processo.

\item Perché gli effetti di un’invenzione possono richiedere tempo?

\item Distingui professione e compito nel caso dell’automazione.

\item Ricostruisci il calcolo energetico dell’ultima parte.

\end{enumerate}
\lessico

\voce{produttività}{prodotto ottenuto rispetto alle risorse impiegate}{إنتاجية}

\voce{innovazione}{introduzione di prodotti o processi migliorati}{ابتكار}

\voce{adozione}{inizio dell’utilizzo di una soluzione}{تبنّي}

\voce{diffusione}{estensione dell’utilizzo}{انتشار}

\voce{automazione}{svolgimento tecnico di compiti prima umani}{أتمتة}

\voce{infrastruttura}{rete o struttura a sostegno di molte attività}{بنية تحتية}

\voce{efficienza}{risultato ottenuto in rapporto alle risorse}{كفاءة}

\lingua
Scrivi quattro relazioni con «a parità di», «per unità di», «nel complesso» e «nel lungo periodo». Indica quali quantità mantieni costanti.
\scrittura
In 280--320 parole descrivi un’innovazione possibile in un’officina o a scuola: problema iniziale, cambiamento, risorse richieste, benefici e limiti.

\end{lettura}
% FINE LETTURA 12

% =================================================================
% INIZIO LETTURA 13 - circa 2320 parole di testo

\begin{lettura}{1}{Economia, ambiente e sviluppo sostenibile}


\subsection*{L'economia dentro l'ambiente}
Ogni attività economica utilizza elementi dell'ambiente. L'agricoltura richiede terra e acqua; l'industria usa materiali ed energia; i servizi dipendono da edifici, trasporti e strumenti. Anche un'attività digitale utilizza dispositivi e infrastrutture fisiche. L'economia non opera quindi in uno spazio separato dalla natura: trasforma risorse e produce conseguenze sull'ambiente in cui viviamo.

L'ambiente offre inoltre condizioni che non sempre vengono vendute attraverso un mercato. Aria respirabile, stabilità degli ecosistemi e disponibilità di acqua hanno valore per le persone e per la produzione. Il fatto che non esista un prezzo diretto per ogni beneficio non significa che quel beneficio sia privo di valore.

Lo \textbf{sviluppo sostenibile} cerca di collegare miglioramento delle condizioni di vita, equità e conservazione delle possibilità future. La parola «sostenibile» richiede quindi una prospettiva temporale. Una scelta può offrire un vantaggio immediato e creare costi successivi. Per valutarla dobbiamo osservare chi riceve il beneficio, chi sostiene il costo e quando questi effetti si manifestano.

\subsection*{Dal costo privato al costo collettivo}
Riprendiamo l'esempio del traffico. Chi utilizza un'automobile considera carburante, tempo e manutenzione. Il viaggio può però contribuire anche a rumore, congestione ed emissioni. Una parte di queste conseguenze ricade su persone esterne alla decisione. Quando non è pienamente inclusa nei costi considerati da chi decide, incontriamo un'esternalità.

Il \textbf{costo privato} è sostenuto dal soggetto che compie l'attività. Il \textbf{costo sociale} considera anche effetti sugli altri. La differenza non significa che chi viaggia desideri danneggiare qualcuno. Significa che le condizioni della decisione possono non rendere visibili tutti gli effetti collettivi.

Immaginiamo che un mezzo produca molto rumore in una strada residenziale. Il conducente può non pagare direttamente per il disturbo provocato. Le famiglie vicine, però, possono perdere tranquillità e riposo. Un'analisi economica deve riconoscere anche questa conseguenza, pur quando è difficile esprimerla con un unico numero monetario.

\subsection*{Inquinamento locale e cambiamento climatico}
L'\textbf{inquinamento atmosferico locale} e il \textbf{cambiamento climatico} sono collegati a varie attività, ma non sono lo stesso fenomeno. Il primo riguarda sostanze presenti nell'aria e conseguenze spesso concentrate in determinate aree. Il secondo riguarda cambiamenti del sistema climatico su periodi lunghi, con effetti distribuiti su scala globale.

La combustione di fonti fossili può contribuire a entrambi, attraverso sostanze e meccanismi differenti. Ridurre alcune combustioni può quindi offrire benefici sia per l'aria locale sia per il clima. Tuttavia, misurare un inquinante in una strada non equivale a misurare direttamente il riscaldamento globale.

Anche \textbf{tempo meteorologico} e \textbf{clima} vanno distinti. Il tempo descrive condizioni di breve periodo, come la temperatura di domani. Il clima riguarda caratteristiche statistiche osservate su periodi lunghi. Una giornata fredda non smentisce da sola una tendenza climatica, così come una sola giornata molto calda non basta a misurarla.

\subsection*{Un dato climatico da leggere bene}
Nel rapporto di sintesi del 2023, l'IPCC indica che la temperatura superficiale globale nel periodo 2011--2020 è stata circa 1,1 gradi Celsius più alta rispetto al periodo 1850--1900. Il rapporto attribuisce il riscaldamento globale alle attività umane, principalmente attraverso le emissioni di gas serra. Si tratta di una sintesi scientifica riferita a periodi precisi.\fonte{clima}

Il dato non significa che ogni città si sia riscaldata esattamente di 1,1 gradi. È una media globale. Non è neppure una previsione per la temperatura di una particolare giornata al Cairo. Le parole «globale», «periodo» e «rispetto a» sono necessarie per interpretarlo correttamente.

Per leggere un numero climatico dobbiamo quindi chiedere quale variabile viene misurata, quale area viene considerata e quale base viene utilizzata. Sono domande simili a quelle poste per il PIL e l'inflazione. La competenza nel leggere i dati attraversa argomenti diversi.

\subsection*{Emissioni e accumulo}
Le \textbf{emissioni} descrivono sostanze rilasciate in un periodo. La \textbf{concentrazione} descrive invece la quantità presente in un determinato volume o insieme. La distinzione richiama quella tra flusso e stock incontrata con reddito e patrimonio. Per alcuni gas, gli effetti dipendono anche dall'accumulo nel tempo.

Un esempio semplice può aiutare, senza rappresentare tutti i dettagli fisici. In un contenitore entra acqua attraverso un rubinetto ed esce da uno scarico. Ridurre l'acqua in entrata non significa necessariamente che il livello cominci subito a scendere: dipende anche da quanta ne esce. Allo stesso modo, rallentare un flusso e ridurre uno stock sono obiettivi differenti.

L'esempio del contenitore è un'analogia, non un modello completo dell'atmosfera. Un buon lettore riconosce sia l'utilità sia i limiti di un paragone. L'analogia aiuta a capire una relazione, ma non sostituisce le misure e le conoscenze scientifiche necessarie per descrivere il fenomeno reale.

\subsection*{Mitigazione e adattamento}
La \textbf{mitigazione} comprende azioni che riducono le cause del cambiamento climatico, per esempio diminuendo emissioni o aumentando assorbimenti. L'\textbf{adattamento} cerca invece di ridurre l'esposizione e la vulnerabilità agli effetti climatici. Sono due orientamenti complementari, non due alternative che si escludono.

Migliorare l'efficienza energetica di un edificio può ridurre l'energia necessaria. Organizzare spazi ombreggiati e procedure per affrontare temperature elevate può ridurre alcuni rischi del caldo. Un intervento può contribuire a entrambi gli obiettivi, ma dobbiamo spiegare attraverso quale meccanismo.

Consideriamo una scuola immaginaria. Una migliore ventilazione e una protezione dal sole possono rendere gli ambienti più utilizzabili. Se permettono anche di ridurre il consumo di elettricità, producono un ulteriore beneficio. Per valutare il progetto occorre però conoscere l'edificio, il clima locale, i materiali e le modalità di utilizzo.

\subsection*{Energia: fonti, potenza e produzione}
L'energia viene ottenuta da fonti diverse e trasformata per gli usi finali. Le fonti fossili comprendono carbone, petrolio e gas naturale. Tra le fonti rinnovabili troviamo sole, vento e acqua in movimento. Il carattere rinnovabile riguarda la disponibilità della fonte, ma non significa che gli impianti siano privi di materiali, costi o conseguenze ambientali.

Dobbiamo distinguere \textbf{potenza} ed \textbf{energia}. La potenza descrive un ritmo di produzione o consumo; l'energia descrive una quantità accumulata durante un intervallo. Un impianto da un megawatt che funzionasse a quella potenza per un'ora produrrebbe un megawattora. Non è corretto usare le due unità come sinonimi.

La potenza installata non indica da sola la produzione annuale. Contano ore di funzionamento, disponibilità della fonte, manutenzione e condizioni della rete. Un impianto solare non produce alla propria potenza massima durante tutte le ore dell'anno. Per confrontare impianti dobbiamo quindi conoscere l'indicatore utilizzato.

\subsection*{Un caso egiziano: Benban}
Il parco solare di Benban, nell'area di Assuan, offre un esempio concreto di investimento energetico in Egitto. Una comunicazione della Banca europea per la ricostruzione e lo sviluppo del gennaio 2019 indicava una capacità contrattualizzata del parco pari a 1.465 megawatt. Il dato descrive un riferimento progettuale di quel momento, non una misura della produzione annuale né un aggiornamento in tempo reale.\fonte{benban}

Per il nostro percorso il caso permette di collegare risorsa naturale, infrastruttura e finanziamento. La disponibilità di sole è una condizione utile, ma un progetto richiede anche terreno, impianti, collegamenti alla rete, manutenzione e accordi economici. Un paesaggio soleggiato non si trasforma automaticamente in elettricità utilizzabile.

Possiamo leggere il numero di megawatt senza attribuirgli significati ulteriori. Non indica quanti megawattora siano stati prodotti in un anno, né quanti lavoratori siano occupati stabilmente. Per rispondere a queste domande servono dati diversi. Il caso mostra perché ogni indicatore deve rimanere collegato alla propria definizione.

\subsection*{Efficienza e consumi totali}
L'\textbf{efficienza energetica} riguarda il servizio ottenuto in rapporto all'energia utilizzata. Una lampada può fornire una certa illuminazione consumando meno elettricità. Questo miglioramento riduce il consumo a parità di servizio, ma il consumo complessivo dipende anche da quante lampade usiamo e per quante ore.

Se una tecnologia rende un servizio meno costoso, le persone possono utilizzarlo di più. Questo possibile aumento dell'utilizzo viene spesso discusso come effetto di rimbalzo. Non significa che ogni miglioramento sia inutile. Significa che il risparmio totale può essere diverso da quello calcolato mantenendo invariati tutti i comportamenti.

Per questo una valutazione deve dichiarare le ipotesi. Se scriviamo «il nuovo sistema riduce i consumi del venti per cento», chiediamo rispetto a quale sistema, con quale utilizzo e in quali condizioni. Una percentuale priva di riferimento non permette di capire il beneficio effettivo.

\subsection*{Regole e incentivi}
Le politiche ambientali possono utilizzare \textbf{regole}, come limiti o requisiti tecnici. Possono anche usare incentivi economici: rendere più costose alcune attività oppure sostenere alternative meno dannose. Gli strumenti modificano condizioni e convenienze, ma richiedono controlli e capacità di applicazione.

Un'imposta su un'attività inquinante può cercare di far considerare a chi decide una parte del costo collettivo. Un contributo a un investimento può ridurre il costo iniziale di una soluzione. Il primo strumento raccoglie entrate; il secondo richiede risorse. Entrambi possono essere progettati in modi differenti.

Non basta scegliere un nome per garantire un risultato. Un incentivo può essere utilizzato anche da chi avrebbe realizzato comunque l'investimento. Una regola può essere difficile da controllare. Perciò bisogna considerare destinatari, comportamenti e possibilità di verifica, oltre all'obiettivo dichiarato.

\subsection*{Un esempio di scelta pubblica}
Immaginiamo una città che voglia ridurre traffico e inquinamento. Può migliorare il trasporto pubblico, creare percorsi pedonali oppure modificare condizioni di accesso e parcheggio. Ogni misura agisce su un passaggio diverso: alternative disponibili, sicurezza, tempo e costo dello spostamento.

Se aumentano i costi dell'automobile senza offrire alternative praticabili, alcune persone possono incontrare difficoltà maggiori di altre. Chi lavora lontano o in orari particolari può avere meno possibilità di adattarsi. Una politica ambientale deve quindi osservare anche la distribuzione degli effetti.

Migliorare il trasporto pubblico richiede investimenti, manutenzione e organizzazione. Acquistare veicoli non basta se le linee non collegano i luoghi necessari o gli orari sono inaffidabili. L'esempio riprende una lezione generale: il bene materiale e il servizio effettivamente offerto non coincidono automaticamente.

\subsection*{Chi sostiene i costi della transizione}
La \textbf{transizione} è un processo di cambiamento tra sistemi di produzione e consumo. Può creare nuove attività e ridurre la domanda di altre. Le persone coinvolte possono avere bisogno di formazione, tempo e sostegno per adattarsi. Il beneficio collettivo di lungo periodo non elimina i costi concentrati nel breve periodo.

Consideriamo un lavoratore specializzato in una tecnologia che viene utilizzata meno. Le competenze possedute possono essere in parte trasferibili, ma non sempre immediatamente. Dire che nasceranno nuovi lavori non dimostra che quella persona possa accedervi senza ostacoli. Contano luogo, requisiti e tempi.

Una discussione sulla \textbf{transizione equa} considera questa distribuzione di opportunità e costi. Non significa bloccare ogni cambiamento. Significa chiedere come renderlo praticabile per i gruppi che affrontano maggiori difficoltà e come utilizzare le risorse disponibili in modo coerente con gli obiettivi.

\subsection*{Materiali e ciclo di vita}
Per valutare un prodotto possiamo osservare il suo \textbf{ciclo di vita}: estrazione dei materiali, produzione, trasporto, utilizzo e fine dell'utilizzo. Un oggetto può consumare poca energia durante l'uso ma richiedere risorse nella produzione. Concentrarsi su una sola fase può quindi offrire un'immagine incompleta.

Riparare, riutilizzare e riciclare sono operazioni differenti. Riparare prolunga la funzione di un oggetto; riutilizzare gli attribuisce un nuovo impiego; riciclare recupera materiali attraverso un processo. Anche il riciclo richiede energia e organizzazione. Non trasforma automaticamente tutti i rifiuti in nuove risorse senza perdite.

Un'impresa può progettare prodotti più facili da mantenere, mentre una città può organizzare raccolte separate. Le scelte individuali funzionano meglio quando esistono infrastrutture e informazioni. Chiedere alle persone di adottare un comportamento senza rendere disponibili i mezzi necessari può limitarne l'efficacia.

\subsection*{Misurare e confrontare}
Un indicatore ambientale può riguardare un totale, un valore per abitante oppure un valore per unità prodotta. Un paese popoloso può avere emissioni totali elevate e un valore pro capite diverso da quello di un paese piccolo. Un'impresa può ridurre le emissioni per pezzo e aumentare quelle totali se produce molto di più.

Nessuna misura è sufficiente per ogni domanda. Se studiamo l'efficienza di una tecnica, il valore per unità può essere utile. Se studiamo la pressione complessiva sull'ambiente, dobbiamo osservare anche il totale. Se discutiamo la distribuzione tra persone, il riferimento pro capite può aggiungere informazioni.

Dobbiamo inoltre indicare il periodo. Una riduzione rispetto a un anno eccezionale può essere diversa da una tendenza di lungo periodo. Le fonti devono specificare metodologia e unità. Il compito del lettore è controllare che il numero utilizzato sostenga davvero la frase che lo accompagna.

\subsection*{Cooperare oltre i confini}
Alcuni effetti ambientali attraversano i confini. Una scelta nazionale può produrre benefici o costi per altri paesi. Nel caso del clima, le decisioni di molti soggetti contribuiscono a un risultato comune. Ritroviamo il problema della cooperazione e degli incentivi a lasciare agli altri il costo dell'intervento.

Gli accordi internazionali cercano di coordinare impegni, informazioni e risorse. La difficoltà riguarda anche differenze nelle capacità economiche, nelle responsabilità e nei bisogni di sviluppo. Un accordo deve trasformare obiettivi generali in azioni, sistemi di verifica e strumenti finanziari.

La prossima lettura approfondirà queste forme di cooperazione. Per ora osserviamo il collegamento: ambiente, economia e istituzioni non sono temi separati. Un progetto tecnico richiede risorse e regole; una regola richiede applicazione; una cooperazione richiede impegni comprensibili e verificabili.

\subsection*{Un controllo prima e dopo}
Per capire se una scuola riduce i consumi possiamo confrontare periodi precedenti e successivi a un intervento. Dobbiamo però considerare anche temperatura esterna, numero di studenti e ore di utilizzo. Un consumo più basso durante una settimana di chiusura non dimostra l'efficacia della nuova lampada o del nuovo impianto.

Il confronto richiede condizioni il più possibile coerenti e una registrazione delle differenze. Possiamo usare consumi per ora di apertura oppure annotare i cambiamenti importanti. Nessun indicatore elimina ogni difficoltà, ma una misura ben descritta permette di discutere il risultato.

Questo metodo collega scienze, matematica e lingua: misuriamo una quantità, costruiamo un confronto e scriviamo una spiegazione. La frase finale dovrebbe indicare ciò che abbiamo osservato e ciò che non possiamo ancora attribuire con certezza all'intervento.

\subsection*{Una domanda finale sullo sviluppo}
Per discutere uno sviluppo sostenibile possiamo partire da un bisogno concreto: spostarsi, abitare, produrre o studiare. Poi confrontiamo le soluzioni disponibili considerando risorse, costi, benefici e conseguenze nel tempo. In questo modo evitiamo di trattare la sostenibilità come una parola positiva priva di contenuto.

Un progetto non deve essere descritto soltanto attraverso ciò che promette. Dobbiamo distinguere un obiettivo, una previsione e un risultato misurato. Possiamo apprezzare una possibilità e riconoscere nello stesso tempo le condizioni necessarie perché si realizzi. È lo stesso metodo utilizzato per le riforme economiche e per l'innovazione.

La lettura dell'ambiente ci chiede quindi di collegare scale diverse: una scelta quotidiana, un'organizzazione produttiva, una politica pubblica e un accordo internazionale. Comprendere questi collegamenti permette di utilizzare il vocabolario economico per descrivere problemi reali senza ridurli a una sola causa o a una sola soluzione.



\begin{center}\begin{minipage}{\linewidth}\centering\small
\begin{tabularx}{\linewidth}{lX}\toprule
Dato & Come va letto\\\midrule
Circa +1,1 $^{\circ}$C & Media globale 2011--2020 rispetto a 1850--1900; IPCC, 2023.\\
1.465 MW & Capacità contrattualizzata di Benban riportata da EBRD nel gennaio 2019; non energia annua.\\\bottomrule
\end{tabularx}
\captionof{table}{Due dati reali con periodi e unità espliciti. Fonti: \ref{src:clima} e \ref{src:benban}.}
\end{minipage}\end{center}


\attivita
\begin{enumerate}

\item Distingui costo privato e costo sociale.

\item Quali periodo e area descrive il dato di 1,1 gradi?

\item Spiega la differenza tra emissioni e concentrazione attraverso l’analogia.

\item Distingui mitigazione e adattamento.

\item Perché i 1.465 MW citati per Benban non sono una produzione annuale?

\item Perché una politica ambientale può avere effetti distributivi?

\item Distingui riparazione, riutilizzo e riciclo.

\end{enumerate}
\lessico

\voce{sostenibile}{che considera condizioni presenti e possibilità future}{مستدام}

\voce{emissioni}{sostanze rilasciate in un periodo}{انبعاثات}

\voce{concentrazione}{quantità presente in un volume o insieme}{تركيز}

\voce{mitigazione}{riduzione delle cause del cambiamento climatico}{تخفيف}

\voce{adattamento}{riduzione della vulnerabilità agli effetti}{تكيّف}

\voce{transizione}{processo di passaggio tra sistemi}{انتقال}

\voce{ciclo di vita}{fasi dalla produzione al fine utilizzo}{دورة الحياة}

\lingua
Individua tre analogie o esempi. Per ciascuno scrivi «Questo esempio aiuta a capire…, ma non dimostra…».
\scrittura
Progetta una proposta di 300--350 parole per ridurre consumi o sprechi a scuola. Specifica bisogno, intervento, risorse, persone coinvolte e indicatore da osservare.

\end{lettura}
% FINE LETTURA 13

% =================================================================
% INIZIO LETTURA 14 - circa 2568 parole di testo

\begin{lettura}{1}{Perché i paesi cooperano?}


\subsection*{Relazioni che attraversano i confini}
Un oggetto utilizzato a scuola può collegare molti paesi. Le materie prime possono provenire da un territorio, i componenti da un altro e il progetto da un altro ancora. Il prodotto viene trasportato, assicurato e venduto attraverso servizi. Anche quando acquistiamo vicino a casa, possiamo partecipare indirettamente a relazioni economiche internazionali.

I paesi cooperano per ragioni diverse: scambiare prodotti, costruire infrastrutture, affrontare rischi comuni oppure organizzare rapporti politici. La \textbf{cooperazione} consiste nel coordinare azioni tra soggetti che mantengono interessi e responsabilità. Non richiede necessariamente che tutti desiderino esattamente la stessa cosa. Richiede però un accordo sulle azioni da compiere e sulle condizioni della collaborazione.

Per leggere un accordo dobbiamo distinguere il problema iniziale, gli obiettivi, gli strumenti e le istituzioni coinvolte. La parola «unione» può indicare realtà molto diverse. Un accordo commerciale, un mercato comune e uno Stato unitario non sono la stessa forma di organizzazione. Questa distinzione guiderà il confronto tra esempi generali, Repubblica Araba Unita e Unione Europea.

\subsection*{Importare ed esportare}
Un'\textbf{esportazione} è una vendita di beni o servizi a soggetti esteri; un'\textbf{importazione} è un acquisto dall'estero. Le due parole descrivono la stessa operazione da punti di vista differenti. Se un'impresa egiziana vende un prodotto a un cliente italiano, l'operazione è un'esportazione per l'Egitto e un'importazione per l'Italia.

Anche i servizi possono attraversare i confini. Una prestazione professionale svolta per un cliente estero può essere un'esportazione di servizi. Nel turismo, un visitatore acquista servizi nel paese in cui soggiorna. Per comprendere il commercio internazionale non dobbiamo quindi pensare soltanto a scatole e container.

Importare non significa necessariamente rinunciare alla produzione interna. Un'impresa può acquistare una macchina estera per produrre beni nel proprio paese. Esportare richiede spesso componenti importati. Le due direzioni possono essere complementari all'interno della stessa filiera. Una frase come «le importazioni sono sempre negative» non descrive questa complessità.

\subsection*{Perché scambiare}
I paesi dispongono di risorse, competenze e condizioni produttive differenti. Lo scambio può permettere di ottenere prodotti difficili o costosi da realizzare localmente, di ampliare la varietà e di raggiungere nuovi clienti. Può anche consentire alle imprese di produrre per un mercato più grande e distribuire alcuni costi su più unità.

Immaginiamo due laboratori. Uno impiega relativamente meno risorse nel produrre tessuti, l'altro nel realizzare strumenti. Una specializzazione parziale e uno scambio possono aumentare le possibilità complessive. L'esempio introduce una logica generale, ma non descrive ogni conseguenza del commercio reale: trasporti, regole, qualità e distribuzione dei benefici contano.

Un vantaggio complessivo non significa che ogni persona benefici nello stesso modo. Alcune imprese trovano nuovi clienti; altre incontrano concorrenti più forti. I lavoratori possono avere bisogno di cambiare attività o formazione. Per questo le politiche commerciali vengono discusse insieme a occupazione, competenze e strumenti di adattamento.

\subsection*{Dipendenza e diversificazione}
Il commercio collega le economie, ma può anche creare \textbf{dipendenze}. Se un'impresa riceve un componente essenziale da un solo fornitore, un'interruzione può bloccare la produzione. Avere più fornitori, scorte o soluzioni tecniche alternative può ridurre alcune vulnerabilità, ma comporta costi.

La \textbf{diversificazione} consiste nel distribuire fonti, attività o relazioni invece di concentrarle in un solo punto. Non elimina ogni rischio. Fornitori diversi possono dipendere dalla stessa rotta o dalla stessa materia prima. Per capire la solidità di una filiera bisogna quindi osservare anche collegamenti meno visibili.

L'obiettivo non è necessariamente eliminare tutti gli scambi. Un paese può cercare di ridurre una dipendenza specifica mantenendo altri rapporti. Autonomia e cooperazione non sono sempre opposti. Collaborare con più partner può essere uno strumento per ampliare le possibilità di scelta.

\subsection*{Ostacoli e regole comuni}
Un \textbf{dazio} è un prelievo applicato alle merci che attraversano una frontiera, in particolare alle importazioni. Può influenzare il prezzo e la convenienza dello scambio. Esistono anche ostacoli non tariffari, legati a procedure, requisiti tecnici o autorizzazioni. Alcuni rispondono a obiettivi di sicurezza o tutela; altri possono rendere difficile l'accesso al mercato.

Immaginiamo un produttore che debba dimostrare in ogni paese, attraverso procedure completamente diverse, che il proprio prodotto rispetta un requisito. Il costo delle verifiche può essere elevato. Regole comuni o accordi di riconoscimento possono semplificare il percorso, purché siano definiti controlli e responsabilità.

Cooperare non significa quindi sempre eliminare regole. Può significare costruire regole condivise, rendere confrontabili i documenti e stabilire procedure per risolvere controversie. La prevedibilità può avere valore economico: un'impresa programma meglio quando conosce le condizioni del mercato in cui intende operare.

\subsection*{Forme diverse di integrazione}
In un accordo di libero scambio i partecipanti riducono determinati ostacoli agli scambi reciproci. Un'unione doganale aggiunge una politica tariffaria comune verso l'esterno. Un mercato comune cerca di facilitare anche il movimento di fattori come lavoro e capitale, oltre agli scambi di beni e servizi.

Un'unione monetaria prevede una moneta condivisa e istituzioni monetarie comuni secondo le regole adottate. Un'unione politica riguarda invece la condivisione o l'unificazione di poteri politici. Queste forme non costituiscono una scala che tutti devono percorrere. Gli accordi reali possono combinare elementi differenti e non avanzano necessariamente nella stessa direzione.

La scelta dipende dagli obiettivi e dalla disponibilità a condividere decisioni. Ridurre un dazio richiede un impegno diverso dal creare una moneta comune o un governo comune. Più ampio è il coordinamento, più diventano importanti rappresentanza, regole decisionali e gestione delle differenze tra partecipanti.

\subsection*{Condividere una decisione}
La \textbf{sovranità} riguarda l'autorità di prendere decisioni su un territorio e una comunità. Attraverso un accordo, gli Stati possono impegnarsi a seguire regole comuni in determinate materie. Il modo in cui questi impegni vengono assunti e applicati varia secondo il tipo di organizzazione.

Immaginiamo due paesi che vogliano costruire un collegamento ferroviario. Devono stabilire il tracciato, la ripartizione dei costi e gli standard tecnici. Devono anche decidere chi gestirà la manutenzione e come affrontare eventuali disaccordi. L'annuncio del progetto è soltanto l'inizio della cooperazione.

Se i benefici arrivano in tempi diversi, può essere difficile mantenere il sostegno. Un paese può sostenere una spesa prima di ricevere entrate. Un accordo può prevedere compensazioni, fasi e verifiche. La collaborazione richiede quindi una struttura capace di gestire cambiamenti e divergenze, non soltanto una dichiarazione di intenzioni.

\subsection*{Un caso storico: Egitto e Siria}
Nel 1958 Egitto e Siria si unirono nella \textbf{Repubblica Araba Unita}, guidata da Nasser. L'unione tra i due paesi terminò nel 1961, quando la Siria si separò. L'Egitto continuò a utilizzare il nome Repubblica Araba Unita fino al 1971. È importante distinguere la durata dell'unione dalla durata del nome.\fonte{rauonu}

Il progetto si inseriva nel contesto del nazionalismo arabo e dell'aspirazione all'unità araba. La condivisione di riferimenti culturali e gli obiettivi politici ebbero un ruolo importante. Non si trattò soltanto di un accordo per ridurre dazi o facilitare vendite: venne costituita un'unione politica.\fonte{raucontesto}

Questo caso permette di comprendere che le ragioni della cooperazione possono essere economiche e politiche insieme. L'aspirazione a una maggiore forza collettiva può accompagnarsi alla ricerca di coordinamento. Tuttavia, per descrivere un progetto storico dobbiamo distinguere le motivazioni dichiarate dalla forma concreta assunta dalle istituzioni.

\subsection*{Differenze dentro un progetto comune}
Egitto e Siria avevano strutture economiche e organizzazioni politiche differenti. Durante l'unione emersero tensioni sul modo in cui erano distribuite le decisioni e sulle trasformazioni economiche. L'accentramento e le diverse posizioni rispetto alle politiche adottate contribuirono alle difficoltà dell'esperienza. Nel settembre 1961 un colpo di Stato in Siria portò alla separazione.\fonte{raudifficolta}

Questa descrizione resta generale. Non attribuisce il risultato a una sola causa e non formula un giudizio sui popoli coinvolti. Una popolazione non ha un unico interesse e un paese comprende gruppi con posizioni diverse. Nella lettura storica dobbiamo evitare di trasformare una decisione di alcune istituzioni in una caratteristica permanente di tutti gli abitanti.

Il caso suggerisce una domanda applicabile anche ad altri accordi: come vengono gestite le differenze tra partecipanti? Condividere un obiettivo non elimina automaticamente differenze nelle risorse, nelle istituzioni e nelle preferenze. Perché una cooperazione duri servono meccanismi che permettano di affrontarle.

\subsection*{Che cosa possiamo confrontare}
Possiamo confrontare esperienze differenti utilizzando alcune dimensioni comuni: obiettivi, partecipanti, istituzioni, decisioni condivise e gestione dei conflitti. Non dobbiamo però supporre che due organizzazioni con la parola «unione» nel nome siano equivalenti. Il confronto serve a riconoscere somiglianze e differenze, non a cancellarle.

Per la Repubblica Araba Unita abbiamo osservato un progetto di unione politica tra due paesi. Per altri casi potremmo osservare un accordo commerciale o una collaborazione tecnica. Il livello di integrazione modifica i problemi organizzativi: coordinare standard di prodotto e governare un territorio comune richiedono istituzioni differenti.

Anche la durata va interpretata con attenzione. Un'esperienza breve può lasciare conseguenze storiche; una collaborazione duratura può incontrare difficoltà e modificarsi. Il numero di anni non sostituisce l'analisi degli obiettivi e del funzionamento.

\subsection*{L'Unione Europea e il mercato unico}
L'Unione Europea offre un diverso esempio di integrazione attraverso trattati e istituzioni comuni. Gli Stati membri mantengono la propria esistenza statale e condividono decisioni in ambiti definiti. Il percorso non è equivalente alla Repubblica Araba Unita e non nasce dallo stesso contesto storico.

Una componente economica centrale è il \textbf{mercato unico}, associato alla libera circolazione di beni, servizi, persone e capitali, secondo le regole previste. Il mercato unico è stato avviato nel 1993 e richiede anche norme comuni e coordinamento. La libertà di movimento non significa assenza di ogni requisito o controllo.\fonte{ue}

Per un'impresa, un mercato più ampio può offrire più clienti e fornitori. Per le persone, la possibilità di spostarsi può ampliare opportunità. Gli effetti dipendono però da competenze, informazioni e condizioni concrete. Anche dentro uno spazio integrato restano differenze territoriali e sociali.

\subsection*{Moneta comune e politiche differenti}
Una parte degli Stati dell'Unione Europea utilizza l'euro. L'area euro e l'Unione Europea non sono quindi due espressioni perfettamente intercambiabili. La politica monetaria comune riguarda l'area euro, mentre molte decisioni di spesa pubblica e tassazione restano nazionali, nel quadro delle regole applicabili.\fonte{bcearea}

Una moneta comune elimina la necessità di cambiare moneta tra i suoi utilizzatori e rende più immediato il confronto di alcuni prezzi. Allo stesso tempo, i partecipanti non possono decidere individualmente un proprio tasso di cambio interno alla moneta comune. La condivisione di uno strumento comporta vantaggi possibili e necessità di coordinamento.

Immaginiamo territori che attraversino fasi economiche differenti. Una politica comune deve considerare un insieme più ampio della singola economia locale. Questo esempio aiuta a capire perché l'integrazione monetaria richieda istituzioni e discussioni su regole e responsabilità. Non basta sostituire il nome stampato sulle banconote.

\subsection*{Infrastrutture e differenze territoriali}
La cooperazione europea comprende anche interventi che mirano a ridurre differenze di sviluppo e sostenere progetti comuni.\fonte{coesione} Per comprendere il meccanismo possiamo immaginare risorse destinate a trasporti, formazione o ricerca. Un progetto deve avere obiettivi, finanziamenti e modalità di realizzazione; la disponibilità di fondi non coincide automaticamente con un risultato.

La logica è simile a quella già incontrata per gli investimenti pubblici. Un'infrastruttura può collegare territori e rendere accessibili attività. La formazione può ampliare competenze. Per valutarne gli effetti bisogna osservare utilizzo, qualità e durata, oltre alla somma spesa.

In un'organizzazione con più membri emerge inoltre la domanda su come distribuire le risorse. I contributi e i benefici non assumono sempre la stessa forma o lo stesso momento. Una strada costruita in un territorio può essere utilizzata anche da imprese di altri territori. Il beneficio economico può quindi attraversare il luogo in cui viene registrata la spesa.

\subsection*{Le istituzioni economiche internazionali}
La cooperazione non avviene soltanto attraverso unioni regionali. Il \textbf{Fondo monetario internazionale}, o FMI, si occupa di cooperazione monetaria, analisi e sostegno finanziario ai paesi secondo strumenti e condizioni specifiche. Non è un governo mondiale e non decide ogni politica economica dei suoi membri.\fonte{fmi}

La \textbf{Banca mondiale} sostiene lo sviluppo attraverso finanziamenti e conoscenze, con strumenti che comprendono prestiti, crediti e contributi secondo le istituzioni e i programmi. Un finanziamento per un progetto deve essere distinto da un trasferimento senza obbligo di rimborso. Leggere lo strumento è necessario per capire l'impegno.\fonte{bm}

L'\textbf{Organizzazione mondiale del commercio}, o OMC, riguarda regole e accordi sugli scambi tra membri. Gli accordi vengono negoziati dagli Stati e dagli altri membri dell'organizzazione secondo le procedure previste. L'OMC non è un'impresa che compra e vende merci per tutti.\fonte{omc}

Queste istituzioni hanno funzioni differenti. Le sigle non devono essere imparate come nomi isolati: vanno collegate a problemi e strumenti. Stabilità monetaria, finanziamento dello sviluppo e regole commerciali sono ambiti connessi, ma non identici. Un paese può partecipare contemporaneamente a più forme di cooperazione.

\subsection*{Leggere un finanziamento internazionale}
Quando leggiamo che un paese «riceve risorse», dobbiamo capire se si tratta di prestito, contributo, investimento o pagamento per un servizio. Un prestito crea un obbligo di rimborso; un investimento può attribuire diritti di proprietà; un contributo segue condizioni diverse. La parola «aiuto» può essere troppo generica.

Dobbiamo inoltre distinguere una somma annunciata, una somma approvata e una somma effettivamente erogata. Un progetto può essere distribuito in fasi. Le erogazioni possono dipendere da verifiche o condizioni. Confondere queste fasi può farci descrivere come già realizzato qualcosa che è ancora previsto.

Anche il \textbf{beneficiario} va identificato. Le risorse possono essere destinate a un'amministrazione, a un'impresa o a un programma. Per comprendere gli effetti sulle famiglie occorre seguire i passaggi: quali attività vengono finanziate, chi le realizza e attraverso quali servizi o redditi raggiungono le persone.

\subsection*{Beni comuni e problemi globali}
Alcuni problemi richiedono una cooperazione ampia perché gli effetti attraversano i confini. Abbiamo incontrato il clima; possiamo pensare anche a reti di trasporto e scambio di informazioni. Ogni paese può beneficiare del contributo degli altri e, nello stesso tempo, avere un incentivo a ridurre il proprio costo.

Gli accordi cercano di rendere gli impegni più chiari e verificabili. Possono prevedere obiettivi, calendari, informazioni condivise e sostegni per chi dispone di minori risorse. La verifica è importante perché una promessa e un'azione non sono la stessa cosa. Anche la fiducia dipende dalla possibilità di osservare ciò che viene fatto.

Le differenze tra paesi non scompaiono. Alcuni hanno maggiori capacità finanziarie, altri bisogni più urgenti o maggiore esposizione ai rischi. Per costruire un accordo occorre discutere come ripartire responsabilità e sostegno. La cooperazione è quindi un processo di negoziazione, applicazione e adattamento.

\subsection*{Un esercizio conclusivo di negoziazione}
Immaginiamo due paesi che vogliano condividere un centro di formazione tecnica. Il primo dispone di edifici e insegnanti; il secondo può contribuire con attrezzature e borse di studio. Entrambi vogliono migliorare le competenze, ma devono decidere quanti studenti ammettere, quali titoli rilasciare e come finanziare la manutenzione.

Un accordo può assegnare compiti e prevedere una verifica annuale. Se le iscrizioni sono inferiori al previsto, i partecipanti devono capire se il problema riguarda informazioni, trasporti o contenuti dei corsi. La cooperazione continua quindi dopo la firma: richiede osservazione e adattamento.

Nel descrivere il caso possiamo riprendere molte parole dell'anno: servizio, investimento, capitale, opportunità, costo e risultato. Possiamo anche distinguere tre frasi: «i paesi desiderano formare tecnici», «i paesi finanziano attrezzature», «gli studenti completano il percorso». La prima esprime un obiettivo, la seconda un'azione e la terza un risultato osservabile. Un testo ben costruito spiega come queste dimensioni sono collegate, senza trattarle come equivalenti. È questa precisione che permette di passare dalla conoscenza delle parole alla comprensione di un ragionamento economico.

\subsection*{Ritrovare il percorso dell'anno}
Siamo partiti da bisogni, beni e servizi. Abbiamo poi osservato produzione, reddito, Stato, credito e riforme. Ora ritroviamo gli stessi concetti nelle relazioni tra paesi. Un accordo internazionale riguarda risorse limitate, scelte, rischi e distribuzione dei benefici, proprio come molte decisioni studiate su scala più piccola.

La scala più ampia rende però importanti nuove istituzioni e nuovi vincoli. Non basta trasferire automaticamente il bilancio di una famiglia al funzionamento di uno Stato o di un'unione. Le analogie aiutano, ma devono essere accompagnate dalle differenze. Una buona lettura conserva entrambe.

Per riassumere un caso di cooperazione possiamo utilizzare cinque domande: chi partecipa, quale problema affronta, che cosa condivide, come decide e come verifica i risultati? La Repubblica Araba Unita, l'Unione Europea e le istituzioni internazionali offrono risposte differenti. Saperle ricostruire significa utilizzare il vocabolario dell'economia per comprendere un testo articolato e discuterlo con precisione.



\begin{center}\begin{minipage}{\linewidth}\centering\small
\begin{tabularx}{\linewidth}{lX}\toprule
Periodo & Repubblica Araba Unita\\\midrule
1958--1961 & Unione politica tra Egitto e Siria.\\
1961--1971 & L’Egitto conserva il nome dopo la separazione siriana.\\\bottomrule
\end{tabularx}
\captionof{table}{Unione e denominazione hanno durate differenti. Fonte: \ref{src:rauonu}.}
\end{minipage}\end{center}


\attivita
\begin{enumerate}

\item Descrivi una stessa operazione come esportazione e importazione.

\item Perché un aumento degli scambi non beneficia necessariamente tutti allo stesso modo?

\item Distingui libero scambio, unione doganale e unione politica.

\item Quando esistette l’unione tra Egitto e Siria? Perché compare anche il 1971?

\item Quali dimensioni permettono un confronto senza equiparare RAU e UE?

\item Distingui FMI, Banca mondiale e OMC.

\item Perché annuncio, approvazione ed erogazione sono fasi diverse?

\end{enumerate}
\lessico

\voce{esportazione}{vendita a soggetti esteri}{تصدير}

\voce{importazione}{acquisto dall’estero}{استيراد}

\voce{dazio}{prelievo sulle merci alla frontiera}{رسم جمركي}

\voce{integrazione}{collegamento attraverso regole e istituzioni comuni}{تكامل}

\voce{sovranità}{autorità di decisione su una comunità e un territorio}{سيادة}

\voce{diversificazione}{distribuzione tra più fonti o attività}{تنويع}

\voce{erogazione}{versamento effettivo di risorse}{صرف الأموال}

\voce{negoziazione}{ricerca di un accordo tra parti}{تفاوض}

\lingua
Scrivi un confronto usando «entrambi», «a differenza di», «mentre» e «in questo caso». Confronta le medesime dimensioni per i due soggetti.
\scrittura
Produci un testo di 350--400 parole: presenta due forme di cooperazione, confrontale con le cinque domande finali e concludi indicando una difficoltà organizzativa comune.

\end{lettura}
% FINE LETTURA 14
\clearpage\section*{Fonti e riferimenti}\addcontentsline{toc}{section}{Fonti e riferimenti}
I riferimenti sostengono definizioni, dati e passaggi storici specifici; esempi immaginari, esercizi e formulazioni didattiche sono originali. Per le fonti storiche si utilizza soltanto il contesto pertinente al periodo descritto. Consultazione: settembre 2026.
\begin{enumerate}[label={[\arabic*]},itemsep=8pt]
\item\label{src:pil} Banca mondiale, glossario del PIL pro capite. Definizione contabile e distinzione tra produzione e popolazione.\par\href{https://databank.worldbank.org/metadataglossary/world-development-indicators/series/NY.GDP.PCAP.CN}{\small Apri la fonte}.
\item\label{src:moneta} Bank of England, How is money created?, 2019. Depositi, prestiti e vincoli bancari.\par\href{https://www.bankofengland.co.uk/explainers/how-is-money-created}{\small Apri la fonte}.
\item\label{src:marx} Karl Marx e Friedrich Engels, \emph{Manifesto del Partito Comunista}, 1848; Karl Marx, \emph{Il capitale}, libro I, 1867. Proprietà dei mezzi di produzione, classi e accumulazione del capitale.\par\href{https://www.marxists.org/archive/marx/works/1848/communist-manifesto/}{\small Apri la fonte}.
\item\label{src:keynes} Fondo monetario internazionale, \emph{What Is Keynesian Economics?} Profilo della teoria di John Maynard Keynes, domanda aggregata e intervento pubblico.\par\href{https://www.imf.org/external/pubs/ft/fandd/2014/09/basics.htm}{\small Apri la fonte}.
\item\label{src:costituzionecinese} Consiglio di Stato della Repubblica Popolare Cinese, \emph{Constitution of the People's Republic of China}, artt. 6, 11, 13 e 15. Proprietà pubblica, settore privato ed economia socialista di mercato.\par\href{https://english.www.gov.cn/archive/lawsregulations/201911/20/content_WS5ed8856ec6d0b3f0e9499913.html}{\small Apri la fonte}.
\item\label{src:inflazione} Banca centrale europea, What is inflation? Definizione, paniere e potere d’acquisto. Esempi numerici del libretto originali.\par\href{https://www.ecb.europa.eu/ecb-and-you/explainers/tell-me-more/html/what_is_inflation.en.html}{\small Apri la fonte}.
\item\label{src:postcovid} Board of Governors of the Federal Reserve System, Jerome H. Powell, \emph{Monetary Policy and Price Stability}, 26 agosto 2022. Ripresa post-pandemica, forte domanda e offerta limitata.\par\href{https://www.federalreserve.gov/newsevents/speech/powell20220826a.htm}{\small Apri la fonte}.
\item\label{src:iran2026} Reuters, \emph{ECB raises interest rates to fight off inflation jump}, 10 settembre 2026. Conflitto iraniano, prezzi dell'energia e inflazione.\par\href{https://www.reuters.com/business/ecb-raises-interest-rates-fight-off-inflation-jump-2026-09-10/}{\small Apri la fonte}.
\item\label{src:bolivia} Juan Antonio Morales e Jeffrey D. Sachs, \emph{Bolivia's Economic Crisis}, NBER Working Paper 2620, 1988. Disavanzo pubblico, espansione monetaria e iperinflazione boliviana.\par\href{https://www.nber.org/papers/w2620}{\small Apri la fonte}.
\item\label{src:bceobiettivo} Banca centrale europea, Measuring inflation and consumer prices. Obiettivo del 2 per cento nel medio termine.\par\href{https://www.ecb.europa.eu/stats/macroeconomic_and_sectoral/hicp/html/index.en.html}{\small Apri la fonte}.
\item\label{src:cbe} Central Bank of Egypt, Monetary Policy Framework. Mandato relativo a prezzi e sistema monetario e bancario.\par\href{https://www.cbe.org.eg/en/monetary-policy/monetary-policy-framework}{\small Apri la fonte}.
\item\label{src:bcemonetaria} Banca centrale europea, Monetary policy. Compiti e strumenti della politica monetaria.\par\href{https://www.ecb.europa.eu/ecb/orga/tasks/monpol/html/index.en.html}{\small Apri la fonte}.
\item\label{src:clima} IPCC, Climate Change 2023: Synthesis Report, Summary for Policymakers, A.1. Riscaldamento globale nel 2011--2020 rispetto al 1850--1900.\par\href{https://www.unep.org/resources/report/climate-change-2023-synthesis-report}{\small Apri la fonte}.
\item\label{src:benban} EBRD, First EBRD funded Egyptian solar plant begins generation, 28 gennaio 2019. Capacità contrattualizzata del parco: 1.465 MW.\par\href{https://www.ebrd.com/home/news-and-events/news/2019/first-ebrd-funded-egyptian-solar-plant-begins-generation.html}{\small Apri la fonte}.
\item\label{src:rauonu} Nazioni Unite, United Arab Republic, voce della Biblioteca digitale. Distinzione tra unione 1958--1961 e nome conservato dall’Egitto fino al 1971.\par\href{https://digitallibrary.un.org/record/1284346?ln=en}{\small Apri la fonte}.
\item\label{src:raucontesto} EBSCO Research Starters, Syria and Egypt Form the United Arab Republic. Contesto dell’unione politica.\par\href{https://www.ebsco.com/research-starters/politics-and-government/syria-and-egypt-form-united-arab-republic}{\small Apri la fonte}.
\item\label{src:raudifficolta} Egypt: A Country Study, Egypt and the Arab World. Unione con la Siria e separazione.\par\href{https://countrystudies.us/egypt/33.htm}{\small Apri la fonte}.
\item\label{src:ue} Consiglio dell’Unione Europea, The EU single market. Quattro libertà e avvio nel 1993.\par\href{https://www.consilium.europa.eu/en/policies/deeper-single-market/}{\small Apri la fonte}.
\item\label{src:bcearea} Banca centrale europea, ECB, ESCB and the Eurosystem. Distinzione tra UE e area euro.\par\href{https://www.ecb.europa.eu/ecb/orga/escb/html/index.en.html}{\small Apri la fonte}.
\item\label{src:fmi} Fondo monetario internazionale, What is the IMF? Funzioni dell’istituzione.\par\href{https://www.imf.org/en/about/factsheets/imf-at-a-glance}{\small Apri la fonte}.
\item\label{src:bm} Banca mondiale, What We Do. Finanziamento e conoscenze per lo sviluppo.\par\href{https://www.worldbank.org/ext/en/what-we-do}{\small Apri la fonte}.
\item\label{src:omc} Organizzazione mondiale del commercio, What is the WTO? Regole e accordi commerciali.\par\href{https://www.wto.org/english/thewto_e/whatis_e/whatis_e.htm}{\small Apri la fonte}.
\item\label{src:coesione} Commissione europea, Cohesion Policy. Riduzione delle differenze territoriali.\par\href{https://ec.europa.eu/regional_policy/policy/what/investment-policy_en}{\small Apri la fonte}.
\end{enumerate}
\end{document}
